\documentclass[11pt]{article}

\usepackage[utf8]{inputenc}
\usepackage[T1]{fontenc}
\usepackage{amsmath,amssymb,amsthm}
\usepackage{mathtools}
\usepackage{geometry}
\geometry{a4paper,margin=1in}
\usepackage{enumitem}
\usepackage{longtable}
\usepackage{array}
\usepackage{booktabs}
\usepackage{hyperref}

\newtheorem{theorem}{Theorem}
\newtheorem{corollary}{Corollary}
\newtheorem{proposition}{Proposition}
\newtheorem{lemma}{Lemma}
\newtheorem{conjecture}{Conjecture}
\newtheorem{obligation}{Obligation}
\theoremstyle{definition}
\newtheorem{definition}{Definition}
\newtheorem{remark}{Remark}
\newtheorem{policyschema}{Policy Schema}

\newcommand{\ObsMatcl}{\mathbf{ObsMatcl}}
\newcommand{\Omat}{\mathcal{O}_{\mathrm{mat}}}
\newcommand{\Minf}{M_\infty}
\newcommand{\plainterms}{\emph{In plain terms.}}
\newcommand{\founding}{\emph{Founding statement.}}

\title{Supplement to A Conditional Construction of Human-Centric Functional Modeling\\
from Globally Lawful Material Flow}
\author{Andy E.\ Williams$^{1,2,*}$\\[4pt]
\small $^1$Nobeah Foundation, Nairobi, Kenya\\
\small $^2$Caribbean Center for Collective Intelligence, St.\ John's, Antigua and Barbuda\\[2pt]
\small $^*$Corresponding author: \texttt{awilliams@nobeahfoundation.org}, \texttt{info@cc4ci.org}}
\date{}

\begin{document}
\maketitle

\begin{abstract}
This supplement provides additional detail for but is not necessary to understand the paper it accompanies. It constructs one object and proves what it generates. Compatible local laws of observable material compositions assemble into one continuous action on a completed material carrier; the action emits canonical interaction generators on every composition; closure and horizon quotienting produce physical functional state objects; a physically realized observer produces registration, canonical conceptual and mathematical cores, declared realized enrichment where needed, and fitness structure; stable refinement produces the global completion; and finite labelled carriers expose the resulting invariants in flat $\mathbb{R}^3$, from which they transport along structure-preserving bridges. This chain is the strong material-generative HCFM theorem, proved here in full, with a coupled-clock witness computed end to end.

Three modules then develop what the completed construction makes measurable, transportable, and controllable. The dimensionality module defines a finite-scale effective registration dimension $d_A(H)$ on a costed operational quotient, proves its invariance and its additivity in the uncoupled product regime, and proves an estimator-calibration theorem: the registration-dimension instrument reports three on a horizon carrying a rank-three rank--metric-regularity certificate, while saturation and exhaustion fix the certified rank. On any declared realization/refinement class whose typical infrared trajectories satisfy those certificates, the generic three-dimensional prediction is globally coherent; causal-order and ultraviolet-profile paths are preregistered empirical realizations. The disciplinary module defines a discipline as a perspective-explicit evaluation-equipped realization, gives the six-clause certified bridge that cross-disciplinary transport must satisfy, specializes invariant transport to that bridge, and formally identifies an observer-inclusion blind spot in general system assessment: when assessment conditions are left outside a system model, a local-prior judgment can require unbounded certified correction before a perspective-explicit global-coherence assessment is reached. It defines the three-valued coherence-budgeted assessment, distinguishes internally executed constraint closure from external certification, proves sufficient conditions for an unbounded blind spot, and proves conditional conditions for a binary limit verdict. It then proves three library results: certified best retained values are monotone under accepted updates; under six named reuse conditions a library holding $n$ compatible modules holds at least $2^n$ semantically distinct certified candidate compositions; and under a coherent task realization with certificate-accounted update work whose stepwise growth is below a factor of two, the guaranteed throughput floor for newly certified real task resolutions accelerates. The alignment module defines operational windows and order transitions, proves that cognitive mass --- signed reach change per unit of registered coupling --- is an order-relative observable invariant with a two-channel population decomposition, and records the nucleation requirement, the two-lever control claim, the realization gate, and the preregistered sign question, each with its falsifier, realization conditions, and positive conditional consequence. Every claim carries its type; each module closes with the ledger that is the single authoritative record of that typing, and ``In plain terms'' glosses restate the load-bearing items for readers without specialist training.
\end{abstract}

\setcounter{theorem}{0}
\setcounter{definition}{0}
\setcounter{proposition}{0}
\setcounter{lemma}{0}
\setcounter{corollary}{0}
\setcounter{remark}{0}
\setcounter{conjecture}{0}
\setcounter{obligation}{0}

\tableofcontents

\section{The construction as one causal object}

A material law completion is the object constructed in this supplement. Its restrictions generate local flows; the flows emit interactions and generate functional models; the models preserve their invariant objects through refinement, representation, and transport.

The complete production chain is
\begin{equation}
\begin{aligned}
\Omat &\xRightarrow{\text{stable compatibility}} (\Minf, \Phi^{\star})
\xRightarrow{\text{restriction and emission}} \{\Phi_C, \mathrm{Emit}(C)\}_C\\
&\xRightarrow{\text{closure / quotient}} \{P_{C,H}\}_{C,H}
\xRightarrow{\text{observer induction}} \{S_{\mathrm{HCFM}}(C, H)\}_{C,H}
\xRightarrow{\text{carrier / transport}} \mathrm{Inv}.
\end{aligned}\tag{S1}
\end{equation}
Here $\Omat$ is a diagram of local material observations. The construction converts it into a global material action exactly when the local law, observation, composition, and refinement conditions commute. This is the sense in which the single-law hypothesis is globally coherent: the construction yields a canonical model that realizes every local observation as a projection of one law. Numbered definitions, theorems, and proofs are the contract throughout; the italic section openers and the ``In plain terms'' glosses restate them for readers with strong general reasoning and no specialist background.

\subsection{Background}

Work in a background in which complete lattices, fixed points of monotone maps, topological inverse limits, quotients by equivalence relations, symmetric monoidal categories, finite presentations, and labelled path categories can be formed. A statement is globally coherent when the named construction exists and its conditions are jointly satisfiable in this background. The result is an existence-and-universality theorem for the stated material-observational conditions.

\begin{remark}[Position in systems literature]
The construction is a conditional systems-theoretic architecture: it treats material composition, observer realization, and invariant transport as typed parts of one production chain. General systems theory supplies the broad concern with system organization [22]; behavioral systems theory supplies an explicit account of systems through admissible trajectories [24]; and compositional applied category theory supplies a mature language for building systems from typed interfaces [23]. The present contribution does not relabel those literatures. It specifies a material-generative HCFM chain in which observer-induced canonical cores, certificate-bearing bridges, and status-typed downstream modules are constructed within one declared assessment discipline.
\end{remark}

\noindent\textbf{Typed symbol register.}

\begin{center}
\begin{tabular}{@{}p{0.16\textwidth}p{0.78\textwidth}@{}}
\toprule
Symbol & Type and role\\
\midrule
$P, A, C, M$ & The physical, registration, conceptual, and mathematical sorts of one HCFM realization.\\
$\Minf, \Phi^{\star}$ & The inverse-limit material completion and its global material action.\\
$P_{C,H}, A_{C,H}$ & Respectively the physical functional-state quotient and the realized operational registration quotient at one composite--horizon pair.\\
$\Pi_K \subset \mathbb{R}^3$ & A labelled inspection carrier for a finite task slice, with an exact decoder.\\
$M^{\mathrm{adm}}_{\mathrm{reg}}, M^{\mathrm{real}}_{\mathrm{reg}}$ & The bare admissible moduli class of registration architectures and its realized subclass (Section 11).\\
$d_A(H)$ & A finite-scale effective dimension evaluated on a declared costed registration datum (Section 11); partial, with $\bot$ as its honest failure value.\\
$D, L$ & Respectively a perspective-explicit disciplinary realization and a persistent certified semantic library (Section 12).\\
\bottomrule
\end{tabular}
\end{center}

\section{Material operands, local observations, and one global law}

Local material observations enter as a directed diagram. The section constructs their global carrier and the one action that has each local law as a projection.

\subsection{Material composition and local law}

\begin{definition}[Material composition object]
A material composition object is a tuple
\[
T = (\mathrm{Mat}, \otimes, |\cdot|, \mathrm{Loc}),
\]
where $\mathrm{Mat}$ is a symmetric monoidal category, $|\cdot| : \mathrm{Mat} \to \mathrm{Top}$ is a strong monoidal carrier functor, and $\mathrm{Loc}_S$ supplies canonical localization maps for every declared factorization of a composite. Thus $|C \otimes D| \cong |C| \times |D|$, and a localized sub-composite flow can be compared on the carrier of its full composite.
\end{definition}

\begin{definition}[Horizon-local lawful material observation]
For a material composite $C$ and horizon $H$, a local lawful observation consists of:
\begin{enumerate}[label=\arabic*.]
\item a nonempty topological material carrier $P^{\mathrm{loc}}_{C,H}$;
\item a continuous reversible action $\Phi_{C,H} : \mathbb{R} \times P^{\mathrm{loc}}_{C,H} \to P^{\mathrm{loc}}_{C,H}$;
\item a horizon $H = (T, \varepsilon, O, Q, \Lambda_H, \mathrm{obs}_H)$; and
\item a monotone inclusive dependency operator $\mathrm{Dep}_{C,H} : \mathcal{P}(P^{\mathrm{loc}}_{C,H}) \to \mathcal{P}(P^{\mathrm{loc}}_{C,H})$.
\end{enumerate}
Inclusiveness means that all flow successors and predecessors accessible inside the horizon belong to the declared dependency set.
\end{definition}

\begin{definition}[Declared stable observation diagram]
\label{def:stable-observation-diagram}
For one declared material construction, let $\Omat : J \to \mathrm{FlowTop}$ be the diagram whose objects are the local lawful observations required by that construction and whose arrows are exactly its declared refinement, calibration, localization, and comparison maps. No arrow between unrelated material systems is implied. The diagram is stably law-indexed when:
\begin{enumerate}[label=(L\arabic*)]
\item[(L1)] each declared refinement arrow carries a continuous surjective carrier map commuting with the local flows;
\item[(L2)] overlapping protocols have calibration maps preserving shared readout classes and labels;
\item[(L3)] every finite compatible family of protocols used in the construction admits a common declared refinement;
\item[(L4$'$)] every cofinal completion subsystem over which an inverse-limit carrier is formed satisfies one declared grounding condition of Lemma~\ref{lem:grounding}; and
\item[(L5)] declared monoidal structural maps and localizations commute with flows and with the comparison maps where composition is used.
\end{enumerate}
Conditions (L1)--(L5) are the formal content of stable observation for the declared diagram. They do not impose a universal index over unrelated composites.
\end{definition}

\begin{definition}[Stable material-observability witness]
\label{def:stable-material-observability-witness}
Let \(S\) be a material subsystem at horizon \(H\). A stable
material-observability witness is a tuple
\[
\mathsf{W}_{S,H}
=
\bigl(
O,C,\mathsf{Prot},\mathsf{out},h,
\mathcal{R},\mathcal{K},\mathcal{D},\mathcal{V},\mathcal{F},\kappa
\bigr),
\qquad C=O\otimes S,
\]
where:
\begin{enumerate}[label=(W\arabic*)]
\item \(O\) is a physically realized material observer;
\item \(\mathsf{Prot}\) is a finite observer-supported probe/readout family
with readout map \(\mathsf{out}\);
\item \(h\) is a realized material history;
\item \(\mathcal{R}=\{(P^{\mathrm{loc}}_j,\Phi_j)\}_{j\in J}\) is the family
of horizon-local lawful material observations induced by the physically used
protocol--horizon occurrences;
\item \(\mathcal{K}\) is the family of physically used refinement,
calibration, localization, and comparison maps;
\item \(\mathcal{D}\) is the dependency, composition, and localization datum;
\item \(\mathcal{V}\) is the observer's declared viable material region;
\item \(\mathcal{F}\) is a declared typed fitness datum: it specifies a
registration continuation-capacity evaluator \(W_A\) and the declared
sort-distinct physical, conceptual, and mathematical fitness objects induced
or realized from it; and
\item \(\kappa\) is either a declared materially realized error-correcting
calibration process or the explicit value \(\bot\), indicating that no strict
adaptive-contraction conclusion is invoked.
\end{enumerate}

The witness is \emph{stable} when:
\begin{enumerate}[label=(S\arabic*)]
\item repeated and overlapping protocols reproduce their operational classes
at the declared resolution;
\item every map in \(\mathcal{K}\) commutes with the corresponding local
flows wherever comparison is used;
\item every finite compatible protocol family used by the witness has a
declared common refinement; and
\item if \(x_j\) is the record of \(h\) at observation node \(j\), then every
comparison map \(a_{ij}\) used by the witness satisfies
\[
a_{ij}(x_j)=x_i.
\]
\end{enumerate}

When \(\kappa\neq\bot\), its contraction certificate consists of a declared
metric \(d_{\mathrm{cal}}\), a domain \(U_{\mathrm{cal}}\), and a constant
\(0\leq\lambda<1\) such that
\[
d_{\mathrm{cal}}\bigl(\kappa(a),\kappa(b)\bigr)
\leq
\lambda d_{\mathrm{cal}}(a,b)
\qquad
\text{for all }a,b\in U_{\mathrm{cal}}.
\]
A material system is materially observable at \(H\) exactly when it carries a
stable material-observability witness.
\end{definition}

\begin{definition}[Observation-closed material-system category]
\label{def:obsmat}
The objects of \(\ObsMatcl\) are the observer-containing material systems
\((C,H)\) for which there exist \(S\), \(O\), and a stable
material-observability witness \(\mathsf{W}_{S,H}\) with \(C=O\otimes S\).
A morphism preserves the declared material flow, admitted protocol labels,
readouts, witness-compatible horizon refinement, and every certificate it
claims to transport.
\end{definition}


\begin{lemma}[Grounding conditions for nonempty completion]
\label{lem:grounding}
Let $(X_j, \rho_{ij})$ be the inverse system on a directed completion subsystem of $J$, with all $X_j$ nonempty. Its inverse limit is nonempty whenever at least one of the following holds: (G1) the subsystem has a greatest calibrated horizon; (G2) every $X_j$ is compact Hausdorff and every $\rho_{ij}$ is continuous; (G3) every $X_j$ is finite; (G4) bonding maps are surjective with finite fibres; (G5) the subsystem has a countable cofinal surjective chain; or (G6) a compatible observation cone is supplied as declared data.
\end{lemma}

\begin{proof}
Under (G1), a point of the greatest coordinate determines the unique compatible thread. Under (G2), the equalizer conditions form a closed family with the finite-intersection property in the compact product, hence have nonempty intersection. Condition (G3) is a discrete instance of (G2). Under (G4), restrict above one index to finite nonempty fibres of the projection to that index and apply (G3). Under (G5), choose successive preimages along the cofinal chain and extend by cofinality. Under (G6), the supplied cone is a point of the limit. These are standard inverse-limit arguments [25, 26, 27].
\end{proof}

\begin{remark}[Why grounding is displayed]
Surjective maps between nonempty carriers do not by themselves guarantee a nonempty inverse limit over an arbitrary directed diagram; the Henkin--Waterhouse phenomenon supplies counterexamples [28, 29]. A greatest horizon is stronger than a merely maximal one. Thus (L4$'$) is a load-bearing condition of the declared completion, not a universal cross-system indexing demand.
\end{remark}

\begin{theorem}[Global material-law completion]
\label{thm:global-material-completion}
Let $\Omat : J \to \mathrm{FlowTop}$ be a declared stably law-indexed observation diagram. Then the inverse limit
\[
\Minf := \varprojlim_{j \in J} P^{\mathrm{loc}}_j
\]
exists and is nonempty. It carries a continuous reversible action
\begin{equation}
\Phi^{\star} : \mathbb{R} \times \Minf \longrightarrow \Minf, \qquad \Phi^{\star}_t((x_j)_j) = (\Phi_j(t)x_j)_j. \tag{S2}
\end{equation}
For each $j \in J$, the projection $p_j$ satisfies $p_j\Phi^{\star}_t = \Phi_j(t)p_j$. The pair $(\Minf, \Phi^{\star})$ is universal among material actions with the same declared local projections, hence unique up to canonical isomorphism.
\end{theorem}

\begin{proof}
Conditions (L1)--(L3) supply the compatible inverse-system data on the declared diagram, and (L4$'$) supplies nonemptiness by Lemma~\ref{lem:grounding}. Define the inverse limit as the subspace of coherent threads of the product. Flow equivariance in (L1) preserves coherence, so (S2) is well defined. The group law, inverses, and continuity are inherited coordinatewise; the projection identities are immediate. Any compatible action cone maps uniquely to the limit by its coordinate maps, and this mediating map is action-equivariant.
\end{proof}

\begin{corollary}[Global coherence of the single-law hypothesis]
A stably law-indexed observable material domain has one canonical global lawful completion. The single-law hypothesis is therefore globally coherent for that domain: every local observation is a lawful projection of one material action, and the completion is universal rather than an arbitrary selection.
\end{corollary}

\noindent\plainterms{} Thousands of photographs of one landscape, from different spots and zoom levels: if every overlapping pair agrees and zooming never contradicts the wide view, there is exactly one landscape they all depict --- and the theorem builds it, together with the one way the whole landscape changes in time. The force of the result is the word one.

\section{Canonical interaction emission from the global law}

The restriction of the global law to a material composite enters as a generator. Localizing all sub-composite generators and inverting their inclusion relation produces the interaction functions emitted by the law.

\begin{definition}[Generator system and localized lifts]
For every observable composite $C = C_1 \otimes \cdots \otimes C_n$, assume the restriction of $\Phi^{\star}$ belongs to a differentiable action class with generator $X_C$. For $S \subseteq [n]$, write $C_S = \bigotimes_{i \in S} C_i$, and let $Y_S$ be the canonical lift of $X_{C_S}$ to $|C|$ determined by $\mathrm{Loc}_S$. Set $Y_\emptyset = 0$.
\end{definition}

\begin{definition}[Emission map]
The emitted interaction component supported on $S \neq \emptyset$ is
\begin{equation}
X_C(S) = \sum_{T \subseteq S} (-1)^{|S|-|T|} Y_T, \qquad \mathrm{Emit}(C) = \{X_C(S) : \emptyset \neq S \subseteq [n]\}. \tag{S3}
\end{equation}
\end{definition}

\begin{theorem}[Canonical interaction-extraction theorem]
\label{thm:canonical-interaction-extraction}
For every observable composite $C$, the emitted family of Definition 10 is the unique family satisfying
\begin{equation}
Y_S = \sum_{T \subseteq S} X_C(T) \quad \text{for every } S \subseteq [n]. \tag{S4}
\end{equation}
In particular, $X_C = \sum_{\emptyset \neq S \subseteq [n]} X_C(S)$. The family is natural under admitted material maps, coherent under refinement of the factorization, and each component has exactly the participant support declared by $S$. When generators do not commute, time-ordered products reconstruct the corresponding flow.
\end{theorem}

\begin{proof}
The Boolean lattice of subsets of $[n]$ has the M\"obius function $\mu(T, S) = (-1)^{|S|-|T|}$. Equation (S3) is M\"obius inversion of the partial-sum relation (S4), so it is unique. Summing (S3) over $T \subseteq S$ cancels all contributions except $Y_S$. Naturality and factorization refinement follow because the lifts $Y_S$ commute with the admitted material maps and localizations; uniqueness carries the equality to the extracted components. The time-ordered reconstruction is the standard reconstruction of a noncommuting sum of generators.
\end{proof}

\begin{definition}[Fundamental emitted interaction]
An emitted component $X_C(S)$ is fundamental at the declared compositional resolution when it is nonzero and no representation of it exists as a sum of emitted components supported on proper subsets of $S$. An empirical interaction realization assigns a component a stable operational signature consisting of participant support, response profile, reconstruction residual, and refinement behavior.
\end{definition}

\noindent\plainterms{} Run each part alone and record what it does; run the combination and subtract everything the parts account for. The remainder is the interaction, and (S3) is inclusion--exclusion doing that subtraction correctly for any number of parts. The theorem's content: the remainder is unique --- interactions are the exact residue of composition, not postulated labels. This is the founding statement's ``single function emits multiple output functions,'' proved.

\section{Closure, histories, and physical functional states}

An emitted and lawful material flow enters this section through a horizon. Dependency closure selects the continuation-complete domain, and observational bisimulation identifies histories whose future labeled behavior cannot be separated at that horizon.

\subsection{Closure}

\begin{definition}[Closure-step operator]
For a candidate domain $D_0 \subseteq P^{\mathrm{loc}}_{C,H}$, define
\begin{equation}
F_{D_0,H}(X) = D_0 \cup X \cup \mathrm{Dep}_{C,H}(X). \tag{S5}
\end{equation}
\end{definition}

\begin{theorem}[Least continuation-complete domain]
\label{thm:least-continuation-complete-domain}
The operator $F_{D_0,H}$ has a least fixed point
\begin{equation}
D^{\sharp}_{C,H} = \mu F_{D_0,H} = \bigcap \{X : D_0 \subseteq X \text{ and } \mathrm{Dep}_{C,H}(X) \subseteq X\}. \tag{S6}
\end{equation}
It is closed under all horizon-bounded forward and backward material evolution.
\end{theorem}

\begin{proof}
Monotonicity of $\mathrm{Dep}_{C,H}$ makes $F_{D_0,H}$ monotone on the complete lattice of subsets. Knaster--Tarski gives the least fixed point and its prefixed-point characterization. Inclusive dependency contains the permitted flow successors and predecessors, so they remain in the fixed point.
\end{proof}

\subsection{Histories and quotient}

\begin{definition}[Admissible history and horizon bisimulation]
An admissible history is a map $h : I_h \to D^{\sharp}_{C,H}$, $I_h \subseteq [-T, T]$, satisfying $h(u) = \Phi_{C,H}(u - v)h(v)$. The horizon labels include all bounded flow translations and any further labels whose enabledness is expressed using material histories, the flow, the closure domain, and the declared horizon readout. A relation is a horizon bisimulation when related histories have equal readout and match each other's labeled transitions forward and backward. Let $\sim_{C,H}$ be the largest such relation.
\end{definition}

\begin{theorem}[Physical functional-state quotient]
\label{thm:physical-functional-quotient}
The greatest horizon bisimulation $\sim_{C,H}$ exists and is an equivalence relation. The quotient
\begin{equation}
P_{C,H} := \mathrm{Hist}_{C,H} / \sim_{C,H} \tag{S7}
\end{equation}
is a well-defined labelled path category. The material flow induces an inverse-closed family
\begin{equation}
r_{C,H}(t)[h] = [T_t h], \qquad |t| \leq T, \tag{S8}
\end{equation}
with the group law on its common bounded domain.
\end{theorem}

\begin{proof}
The bisimulation operator on binary relations is monotone; Knaster--Tarski yields its greatest fixed point. Identity, inverse, and composition of bisimulations yield reflexivity, symmetry, and transitivity. Matching transitions make arrows independent of representatives. Flow translations preserve admissibility by Theorem 14; deterministic flow labels make (S8) well-defined. The action law and inverses descend coordinatewise to equivalence classes.
\end{proof}

\begin{proposition}[Reversibility and adaptive contraction are distinct types]
\label{prop:reversibility-adaptive-contraction}
If an invertible map and its inverse are nonexpansive on one metric space, the map is an isometry. Hence a strictly contractive adaptive correction map belongs to a separate type from the reversible navigation family $r_{C,H}$.
\end{proposition}

\begin{proof}
Nonexpansiveness of the map gives $d(rx, ry) \leq d(x, y)$; nonexpansiveness of its inverse gives the reverse inequality.
\end{proof}

\noindent\plainterms{} Closure packs for the trip: everything that could matter within the time-and-precision budget, nothing that cannot. The quotient then merges any two histories no allowed observation could ever tell apart. What remains is the world at the resolution at which it can make a difference --- the physical functional state object. The small proposition after it records a load-bearing distinction: navigating (reversible) and self-correcting (strictly contracting) are different kinds of motion, and no map can be both.

\section{The four-space HCFM schema}

The physical quotient becomes an HCFM instance when observation is physically realized. The observer supplies the operational distinctions that generate registration and the material conditions that generate viability, fitness, and correction.

\begin{definition}[Physically realized observer]
An observer-containing composite is $C = O \otimes S$ together with a finite probe family supported in $O$, a readout of the observer state at the declared resolution, and material sensitivity: some protocol separates two distinguishable material histories. Two histories are operationally equivalent when no admissible finite protocol of probes and flow segments separates their readouts.
\end{definition}

\begin{definition}[Canonical HCFM realization]
\label{def:canonical-hcfm-realization}
For an observable material system $(C, H)$, its canonical HCFM schema instance is
\begin{equation}
\begin{aligned}
S^{\circ}_{\mathrm{HCFM}}(C, H) = \big(&\{\mathrm{FMA}^{\sigma,\circ}_{C,H}\}_{\sigma \in \{P,A,C,M\}};\ \mathrm{Reg}_{P \to A}, \mathrm{Ext}_{K^{\circ} \to P(A)},\\
&\mathrm{Int}_{K^{\circ}_{\mathrm{def}} \to M^{\circ}_{\mathrm{def}}},\ \mathrm{Real}_{M^{\circ} \to \mathrm{Def}(A(C,H))}\big),
\end{aligned}\tag{S9}
\end{equation}
where $P, A, C, M$ denote the physical, registration, conceptual, and mathematical sorts. For each sort $\sigma$, the corresponding column is
\[
\mathrm{FMA}^{\sigma,\circ}_{C,H} = (\Omega^{\sigma,\circ}, R^{\sigma,\circ}, S^{\sigma,\circ}, W^{\sigma,\circ}),
\]
with a typed path category, inverse-closed declared dynamics, and a sort-distinct fitness object. The $C$- and $M$-columns are the canonical cores of Definitions~\ref{def:canonical-conceptual-core} and \ref{def:canonical-mathematical-core}. A realized HCFM instance may enrich these cores only by declared data equipped with a typed comparison interface and residual conditions.
\end{definition}

\begin{definition}[Protocol/readout context and canonical conceptual core]
\label{def:canonical-conceptual-core}
Let $\mathrm{Prot}_{C,H}$ be the admissible finite protocols of a physically realized observer, and let $\mathrm{out}(\pi, h) \in \mathrm{Val}_\pi$ be their declared finite-resolution readouts. Operational equivalence identifies histories with equal readouts for every $\pi$, giving $A_{C,H}$. Each readout descends uniquely to $\mathrm{out}(\pi, -) : A_{C,H} \to \mathrm{Val}_\pi$. The observational context has objects $A_{C,H}$, attributes $\mathrm{Att} := \{(\pi, v)\}$, and incidence $a \Vdash (\pi, v)$ iff $\mathrm{out}(\pi, a) = v$. Its formal concepts $(E, J)$, defined by the Galois closure $E' = J$, $J' = E$, form the complete lattice $K^{\circ}(C, H)$ [30]. Its order paths are intent-certified extent inclusions, its bounded dynamics are induced by protocol composition with bounded flow translations, and its fitness is $W_{K^{\circ}}(E, J) = \sup_{a \in E \cap V_{C,H}} W_A(a)$.
\end{definition}

\begin{definition}[Canonical mathematical core]
\label{def:canonical-mathematical-core}
The observational signature $\Sigma(C, H)$ has a state sort $A$, declared value sorts, bounded flow symbols $r_t$, and readout symbols $\mathrm{out}_\pi$. The registered structure $A(C, H)$ interprets these on the declared bounded/slack domain. The canonical mathematical core is the syntactic Lindenbaum--Tarski category
\[
M^{\circ}(C, H) := \mathrm{Syn}(\mathrm{Th}(A(C, H))),
\]
of contexts and provable-equivalence classes of definable functional relations [31, 32]. Its certified subcore is generated by the sentences carrying the paper's declared protocol-verification certificates. The full conceptual core remains the complete lattice of Definition~\ref{def:canonical-conceptual-core}; its total semantic relation to registration is the extent map $\mathrm{Ext}(E, J) = E \subseteq A_{C,H}$. Define the finitarily definable interface
\[
K^{\circ}_{\mathrm{def}}(C, H) := \{(E, J) \in K^{\circ}(C, H) : \exists \text{ a finitary } \Sigma(C, H)\text{-formula } \varphi_E(x) \text{ with } E = \varphi_E^{A(C,H)}\}.
\]
The classifier $\mathrm{Int}_{K^{\circ}_{\mathrm{def}} \to M^{\circ}_{\mathrm{def}}}$ sends $(E, J)$ to the provable-equivalence class of any such $\varphi_E$. It is well defined because $\mathrm{Th}(A(C, H))$ is the complete observational theory of $A(C, H)$, so any two formulas with the same extent are equivalent in that theory. Concepts outside $K^{\circ}_{\mathrm{def}}(C, H)$ remain in the full lattice with the typed residual $\rho_{\mathrm{def}}(E, J) = \text{nondefinable}$. The realization bridge $\mathrm{Real}_{M^{\circ} \to \mathrm{Def}(A(C,H))}$ is the canonical semantic interpretation functor into definable contexts and relations over the registered structure. A material implementation of an interpreted object is not supplied by this functor alone: it is declared realization enrichment with a typed material-realization interface and residual conditions.
\end{definition}

\begin{theorem}[Derived observer columns]
\label{thm:derived-observer-columns}
For a physically realized sensitive observer equipped with declared fitness
realization data: (i) operational equivalence yields the registration quotient
and every protocol readout descends to it; (ii) the protocol/readout context
generates the complete canonical conceptual core $K^{\circ}(C, H)$; (iii) the
observational signature and registered structure generate the canonical
mathematical core $M^{\circ}(C, H)$; (iv) the total extent map, the classifier
on the declared finitarily definable interface
$K^{\circ}_{\mathrm{def}}(C, H)$, and the semantic interpretation functor into
$\mathrm{Def}(A(C, H))$ give the displayed bridges, while concepts outside
$K^{\circ}_{\mathrm{def}}(C, H)$ retain their typed non-definability residual;
(v) the declared fitness datum supplies the sort-distinct fitness objects; a
strict adaptive-contraction conclusion is available only when the declared
material calibration process carries a contraction certificate; and (vi) any
structure beyond the canonical conceptual and mathematical cores is declared
realization enrichment with a typed comparison interface.
\end{theorem}

\begin{proof}
Readout descent follows immediately from the definition of operational
equivalence. Formal concept analysis gives the complete concept lattice of
Definition~\ref{def:canonical-conceptual-core}; bounded flow translations act
on attributes by protocol composition and therefore preserve the corresponding
extents on their common domain. The total semantic bridge from that full
lattice is its extent map.
Definition~\ref{def:canonical-mathematical-core} constructs
$K^{\circ}_{\mathrm{def}}(C, H)$ without replacing $K^{\circ}(C, H)$: on that
declared interface, the classifier is well defined by complete
observational-theory equivalence, and the standard semantic interpretation
functor has codomain $\mathrm{Def}(A(C, H))$. The stated non-definability
residual records the complement of that interface. The fitness clause consumes
the declared observer fitness datum, and the contraction clause consumes a
materially realized calibration process together with its certificate;
Proposition~\ref{prop:reversibility-adaptive-contraction} keeps reversible
navigation distinct from adaptive contraction. The enrichment clause is a
type boundary, not an additional existence claim.
\end{proof}

\begin{theorem}[Observability closure]
\label{thm:observability-closure}
Let \(\mathsf{W}_{S,H}\) be a stable material-observability witness. Then:
\begin{enumerate}[label=(\roman*)]
\item its local observations and physically used maps determine a declared
stably law-indexed observation diagram
\[
\mathcal{O}_{\mathsf W}:J_{\mathsf W}\longrightarrow \mathrm{FlowTop};
\]
\item on every completion subsystem used by the witness, the coordinate
records of \(h\) form a compatible cone and therefore satisfy grounding
condition \textup{(G6)} of Lemma~\ref{lem:grounding};
\item \(\mathcal{D}\) supplies the input to dependency closure and horizon
bisimulation, while \(O\), \(\mathsf{Prot}\), \(\mathsf{out}\),
\(\mathcal{V}\), \(\mathcal{F}\), and, where invoked, \(\kappa\), supply the
observer realization data; and
\item \((O\otimes S,H)\in\ObsMatcl\), with canonical HCFM realization
\[
S^{\circ}_{\mathrm{HCFM}}(O\otimes S,H).
\]
\end{enumerate}
The canonical realization includes the registration quotient and canonical
conceptual and mathematical cores, together with the fitness objects declared
by \(\mathcal{F}\). It includes strict adaptive contraction only when
\(\kappa\neq\bot\) carries the displayed contraction certificate.
\end{theorem}

\begin{proof}
Index the witness's local observations by the finite protocol--horizon
occurrences physically used in the construction, and place in
\(\mathcal{O}_{\mathsf W}\) exactly the maps in \(\mathcal{K}\). Stability clauses
\textup{(S1)}--\textup{(S3)} supply the refinement, calibration, and
common-refinement conditions of
Definition~\ref{def:stable-observation-diagram}; \(\mathcal{D}\) supplies
the declared monoidal/localization coherence where composition is used.

For each node \(j\), let \(x_j\) be the record of the same realized material
history \(h\). Stability clause \textup{(S4)} gives
\(a_{ij}(x_j)=x_i\) for each physically used comparison arrow. Hence
\((x_j)_j\) is a compatible cone and provides grounding condition
\textup{(G6)} on every declared completion subsystem.

Theorem~\ref{thm:global-material-completion} therefore supplies the global
material completion. Theorem~\ref{thm:physical-functional-quotient} applies
to the dependency and horizon data, and
Theorem~\ref{thm:derived-observer-columns} constructs the canonical
observer-induced columns. Definition~\ref{def:obsmat} then gives
\((O\otimes S,H)\in\ObsMatcl\).
\end{proof}

\begin{corollary}[Recursive observability closure]
\label{cor:recursive-observability-closure}
Stable material observability is closed under compatible observer-inclusive
composition and grounded horizon refinement:
\begin{enumerate}[label=(\roman*)]
\item compatible finite compositions of witness-bearing observer-containing
systems carry the induced product witness;
\item a finite nested observation chain is a repeated compatible material
composition and therefore remains in \(\ObsMatcl\); and
\item a directed witness-preserving horizon-refinement tower whose record
threads remain compatible has the corresponding canonical completion.
\end{enumerate}
Consequently, every material system observable at its own declared stable
horizon receives the canonical HCFM construction, and this applicability is
closed under the material operations through which observation is extended.
\end{corollary}

\begin{proof}
For compatible composites, combine the material carriers, observer/protocol
families, localized flows, dependency data, fitness data, and record threads.
The resulting data satisfy
Definition~\ref{def:stable-material-observability-witness}. Nested observation
is the finite iteration of this construction. In a directed refinement tower,
witness compatibility provides the record cone required by
Theorem~\ref{thm:global-material-completion}. Theorem~\ref{thm:observability-closure}
then applies at each stage.
\end{proof}

\begin{remark}[Readout-level scope]
\label{rem:readout-level-scope}
The witness is stated at the material-realization level. A stochastic,
irreversible, hybrid, or dissipative description is admitted when it occurs as
a declared observational or readout-level presentation of a witness-bearing
material process. Its inclusion does not require that the readout presentation
itself be reversible; it requires that the physically used material
observation data instantiate the lawful local-flow construction.
\end{remark}

\noindent\plainterms{} Material observability is not a label assigned after a
system has passed some separate membership test. It is a physical witness:
observer, probes, readouts, repeated comparisons, and compatible records of
one material history. Those records supply the observation diagram and its
grounding thread; the observer supplies the quotient and canonical HCFM
columns. Adding observers or refining horizons repeats this same construction.


\begin{theorem}[Variance-aware canonical-core transport]
\label{thm:variance-aware-transport}
A protocol-explicit material morphism induces covariant maps on physical and registration quotients. On canonical conceptual and mathematical cores, it induces contravariant transport on the observational fragment generated by readout-comparison formulas: concept extents pull back by preimage, and observational formulas pull back by signature translation. Quantified pullback requires surjective registration refinement; full-language transport outside the certified observational fragment is a typed partial bridge with a residual. Thus the general morphism of realized HCFM schemas is a typed dynamic-equivariant bridge, not an unrestricted plain functor on all conceptual or mathematical language.
\end{theorem}

\begin{proof}
A material morphism preserves admissible histories, readouts, and operational equivalence, hence induces the physical and registration maps. Preimage preserves intersections of readout-defined extents, giving contravariant meet-preserving concept transport. Readout-comparison formulas pull back along the induced registration map; state-sort quantifiers require surjectivity to lift witnesses. A coarse-graining can make a bounded flow step observationally indistinguishable from identity while the fine system still distinguishes it, so unrestricted forward preservation of all state-sort identities is unavailable. The stated residual records precisely this failure outside the observational fragment.
\end{proof}

\noindent\plainterms{} A physical observer first produces a registration quotient. From that quotient, its own protocol/readout distinctions generate a minimal concept lattice, and the truths of its registered operations generate a minimal mathematical syntax. These are actual constructed columns, not labels added after quotienting. Richer conceptual or mathematical practice may be attached, but it is declared enrichment with an explicit comparison interface. When one observer is coarsened into another, readings pull back reliably while unresolved raw-state identities need not: that is why HCFM transport uses typed bridges with residuals.

\section{Composition closure and canonical-core bridge discipline}

An HCFM instance enters as the image of one observable material system. Compatible material composition produces the instance of the composite, including the interaction terms that are absent from a mere side-by-side juxtaposition.

\noindent\textbf{Recall.}
Definition~\ref{def:obsmat} defines \(\ObsMatcl\) by stable
material-observability witnesses. Theorem~\ref{thm:observability-closure}
shows that this definition is constructive: the witness generates the
observation diagram, its grounding cone, the physical quotient, and the
canonical observer-induced HCFM realization.

\begin{theorem}[HCFM canonical-core bridge construction]
\label{thm:hcfm-canonical-core-bridge}
The assignment $(C, H) \mapsto S^{\circ}_{\mathrm{HCFM}}(C, H)$ sends a material morphism to a typed dynamic-equivariant bridge: its physical and registration components are covariant, while its canonical conceptual and mathematical observational components are the contravariant maps of Theorem~\ref{thm:variance-aware-transport}. The bridge records any coarse-graining, value-map, fragment, or certificate residual. On domains where these residuals vanish, identities and composition are inherited from the material maps and the canonical-core maps compose accordingly.
\end{theorem}

\begin{proof}
A flow- and readout-preserving material map sends admissible histories to admissible histories and preserves operational equivalence, giving the covariant physical and registration components. The canonical conceptual and mathematical components are constructed by Theorem~\ref{thm:variance-aware-transport}; the extent, definability-interface classifier, semantic interpretation, and residual interfaces are those of Definitions~\ref{def:canonical-conceptual-core} and~\ref{def:canonical-mathematical-core}. Thus all components are defined before bridge formation, and the stated identity/composition laws hold on their declared certified domains.
\end{proof}

\begin{theorem}[Composition closure]
\label{thm:composition-closure}
For compatible objects \((C,H_C),(D,H_D)\in\ObsMatcl\), the composite
\(C\otimes D\) lies in \(\ObsMatcl\) with the induced horizon. It admits
declared component-to-composite typed bridges from
\(S^{\circ}_{\mathrm{HCFM}}(C,H_C)\) and
\(S^{\circ}_{\mathrm{HCFM}}(D,H_D)\), together with the nonzero emitted
interaction components indexed across both factors. Each bridge declares the
structures it preserves; any coupling-induced coarse-graining or distinction
loss is recorded as a typed residual. Thus HCFM is closed under compatible
material composition and under the internal composition of its typed
operations.
\end{theorem}

\begin{proof}
The strong monoidal carrier and localization maps provide the composite
material carrier, generator lifts, and the declared component-to-composite
maps. Theorem~\ref{thm:canonical-interaction-extraction} supplies the
cross-composite components. The product observer/protocol construction
provides an operational quotient when the component protocols are compatible.
Theorems~\ref{thm:least-continuation-complete-domain},
\ref{thm:physical-functional-quotient}, and
\ref{thm:derived-observer-columns} then construct the composite HCFM
instance. Path composition is inherited from the quotient categories.
\end{proof}

\noindent\plainterms{} The schema assignment says: the model-building recipe is the same recipe for every witness-bearing system, and translating between systems translates their models through typed bridges. Composition closure says: plug two modeled systems together and the combination is modeled by the same recipe --- carrying declared component bridges plus the genuinely new interaction terms that exist only because the two now touch. Together these establish the projection claim for materially observable systems: one schema because the schema is the construction, not a vocabulary.

\section{Stable horizon completion and observable invariants}

The same material system is now represented at several horizons. Compatible refinement assembles its functional models, and the properties that survive the quotient become the model's own invariant objects.

\begin{definition}[Observable invariant]
For a horizon $H$, a map $q : \mathrm{Hist}_{C,H} \to V$ is an observable invariant when $q(h) = q(k)$ whenever $h \sim_{C,H} k$. It is flow-stable when $q \circ T_t = q$ for all declared flow translations.
\end{definition}

\begin{proposition}[Invariant factorization]
A map $q$ is an observable invariant exactly when there is a unique $\bar{q} : \mathrm{Ob}(P_{C,H}) \to V$ such that
\begin{equation}
q = \bar{q} \circ \pi_H. \tag{S11}
\end{equation}
It is flow-stable exactly when $\bar{q} \circ r_{C,H}(t) = \bar{q}$.
\end{proposition}

\begin{proof}
The quotient map $\pi_H$ identifies exactly the equivalence classes of $\sim_{C,H}$. Class-constancy is therefore equivalent to unique factorization. The flow statement follows from $\pi_H T_t = r_{C,H}(t)\pi_H$.
\end{proof}

\begin{theorem}[Horizon completion]
\label{thm:horizon-completion}
A stably law-indexed compatible horizon family induces a compatible inverse system of physical quotient categories and HCFM realizations. Its inverse limit $P_\infty$ exists, carries the flow threadwise, projects to every horizon model, and is universal among compatible HCFM completion cones.
\end{theorem}

\begin{proof}
Refinement maps preserve admissibility, observation, and labels; thus they descend to quotient functors. The declared horizon-completion diagram inherits the grounding condition required by Definition 4, so Lemma~\ref{lem:grounding} supplies a nonempty compatible carrier where nonemptiness is claimed. The inverse-limit construction and universal property then follow as in Theorem~\ref{thm:global-material-completion}, now in the category of quotient functional models.
\end{proof}

\begin{proposition}[Invariant transport]
Let $F : P_S \to P_T$ be a typed, label-compatible, dynamic-equivariant bridge. Then
\begin{equation}
F^{*} : \mathrm{Inv}(P_T) \to \mathrm{Inv}(P_S), \qquad \bar{q} \mapsto \bar{q} \circ F \tag{S12}
\end{equation}
is well-defined and preserves flow-stability. When $F$ is an equivalence, transport is bijective.
\end{proposition}

\begin{proof}
The composite $\bar{q} \circ F$ is a map on the source quotient, hence an invariant. Equivariance gives $\bar{q}F r_S(t) = \bar{q}r_T(t)F = \bar{q}F$. A quasi-inverse supplies the inverse transport under equivalence.
\end{proof}

\noindent\plainterms{} Equation (S11) is the founding statement's last sentence, proved: an observable property IS a map that factors through the system's own tell-apart quotient --- an invariant object in the system's own functional model --- and (S12) is the transport: any translation that respects the machinery which produced a property carries the property with it.

\section{Flat three-dimensional carriers and human performance}

A finite functional slice enters as a typed incidence presentation. The carrier theorem exposes it in flat $\mathbb{R}^3$ without replacing its functional geometry by Euclidean geometry.

\begin{definition}[Finite task slice and carrier]
A finite task slice is a finite typed atom presentation of states, relations, labels, and incidence data from an HCFM realization. Its simple incidence skeleton replaces each higher-arity relation by a relation node and labelled incidences. For an injective identifier $\nu$ on a countable presentation universe, place an atom $u$ at
\begin{equation}
\iota(u) = (\nu(u), \nu(u)^2, \nu(u)^3) \in \mathbb{R}^3. \tag{S13}
\end{equation}
The labelled straight-line realization is $(\Pi_K, \lambda_K)$.
\end{definition}

\begin{theorem}[Faithful coherent three-dimensional carrier]
\label{thm:faithful-3d-carrier}
Every finite task slice $K$ has a crossing-free labelled carrier $(\Pi_K, \lambda_K)$ in flat $\mathbb{R}^3$ with a decoder that recovers the typed presentation exactly. For $K \subseteq K'$, the carrier of $K'$ restricts to the carrier of $K$. Three dimensions are minimal for uniform straight-line crossing-free realization of the simple incidence-skeleton class.
\end{theorem}

\begin{proof}
No four distinct points of the moment curve in $\mathbb{R}^3$ lie in one plane. Two nonincident straight-line edges can cross only if their four endpoints are coplanar; therefore the simple skeleton is crossing-free. The labels preserve all relation and type information excluded from the bare placement, so decoding is exact. A fixed identifier map gives restriction coherence. Nonplanar members of the skeleton class preclude a uniform two-dimensional straight-line crossing-free realization.
\end{proof}

\begin{definition}[Human performance functional]
A human performance functional is a declared empirical functional $J_H(\Pi, \lambda)$ on faithful labelled carriers, measured for a specified population and task distribution. It may combine decoding accuracy, time, error resilience, learnability, and other declared measures.
\end{definition}

The carrier theorem establishes the human-accessible structural condition of the founding statement: every finite slice is faithfully available in a native flat Euclidean three-dimensional carrier --- proximity, adjacency, and reachability in the drawing are genuine structural properties, not artifacts of an uncontrolled dimensionality reduction --- and the optimization statement $(\Pi^{\star}, \lambda^{\star}) = \arg\max J_H$ is an empirical theorem relative to the declared population and task distribution. The theorem's ``three'' is a property of the inspection carrier; the dimensionality that a realized observer registers is the separate measured quantity constructed in Section 11, and the comparison with brane-world confinement is analogical only, since those models localize physical interaction channels rather than supply representational carriers [5, 6].

\noindent\plainterms{} Whatever the model's own geometry --- curved, tangled, high-dimensional --- any finite piece can be drawn in ordinary flat 3D with straight lines that never cross, plus labels from which a decoder rebuilds the piece exactly. The trick: place points on a curve so twisted that no four ever share a plane, so no two connecting lines can collide. Three dimensions always suffice, and no fewer suffice in general.

\section{Supporting HCFM modules}

The global construction supplies the primary causal picture. These modules preserve the same picture at additional functional resolutions: path geometry measures it, presentation operations expose it, and meta-lifts carry it across finite orders.

\subsection{Intrinsic path geometry and coarse curvature}

\begin{definition}[Cost, identification, and coarse curvature]
For a physical quotient $P_{C,H}$, a declared label cost $c : \Lambda_H \to [0, \infty)$ assigns each labelled path the sum of its edge costs. The infimum over paths from $x$ to $y$ is the directed distance $\vec{d}(x, y)$. Let $d_b(x, y) = \max\{\vec{d}(x, y), \vec{d}(y, x)\}$, and identify states at zero $d_b$-distance. A declared probability measure assignment $m_x$ on the metric quotient defines the coarse curvature
\[
\kappa(x, y) = 1 - \frac{W_1(m_x, m_y)}{d_b(x, y)}
\]
for distinct comparable states, where $W_1$ is the induced transport distance.
\end{definition}

\begin{proposition}[Geometry module]
The path cost of Definition 34 is an extended directed quasi-pseudometric. Symmetrization and zero-distance identification yield a metric quotient. The coarse-curvature expression is well-defined on its declared domain. When the cost is flow-equivariant, every reversible physical flow operator acts by an isometry of the metric quotient.
\end{proposition}

\begin{proof}
Concatenating paths gives the directed triangle inequality, and zero-length paths give reflexivity. Symmetrization is symmetric and satisfies the triangle inequality; quotienting by zero distance supplies separation. The measure and metric quotient provide the domain for $W_1$. Flow equivariance preserves the costs of paths and their inverses, hence preserves the infimum distance.
\end{proof}

\subsection{Presentation operations and target-relative recovery}

\begin{definition}[Presentation and synchronization residue]
A finite typed atom presentation records state, relation, incidence, label, and type atoms of an HCFM slice. The operation classes $G^{\sigma}(R)$ and $L^{\sigma}(R)$ respectively generate local expansion/rewiring and local contraction/identification moves. For a declared extension-stable target $Q$, the synchronization residue is the finite set of target-relevant atoms present in a global synchronization but absent from the accessed local presentation.
\end{definition}

\begin{theorem}[Presentation factorization and recovery]
\label{thm:presentation-factorization}
Every finite typed presentation transition factors into a finite word in the $G/L$ operation classes. If the accessed registration state presents the local synchronization of an extension-stable target, then an empty synchronization residue or a factorization certificate yields recovery of the target from that registration state.
\end{theorem}

\begin{proof}
A finite presentation transition decomposes into primitive creation, deletion, retyping, incidence change, splitting, and identification steps. These are generated by the declared $G/L$ operations. When the residue is empty, extension stability makes the target value independent of atoms outside the local presentation; a factorization certificate records an equivalent finite derivation through the accessed state.
\end{proof}

\subsection{Horizon refinement and finite-order meta-lift}

\begin{definition}[Faithful encoding-and-lift scheme]
A faithful encoding-and-lift scheme assigns each finitely presentable HCFM realization a typed finite presentation, an exact decoder, and a lift operation that preserves the domain types, dynamic bases, path composition, bridges, and fitness objects. It is flow-derived when the lifted physical column is induced by the encoded material flow.
\end{definition}

\begin{theorem}[Refinement and meta-lift closure]
\label{thm:refinement-meta-lift-closure}
A horizon refinement induces a canonical abstraction functor between the corresponding physical quotient categories and, when the realization data commute with refinement, between HCFM instances. Under Definition 38, every finite-order lift preserves the HCFM architecture; flow-derivedness is preserved at each order when the lift is flow-derived.
\end{theorem}

\begin{proof}
Refinement maps preserve observation and labeled transitions, so bisimulation classes map to bisimulation classes; this defines the quotient abstraction functor. The decoder and type-preservation conditions of Definition 38 carry each component of the HCFM tuple and its bridges to the lifted presentation. Induction on lift order gives finite-order closure.
\end{proof}

\noindent\plainterms{} The meta-lift is the machinery of level-jumping: whole compositions at one order become single primitives one order up, with a decoder guaranteeing nothing is lost. The alignment module's ``coherent order transition'' is exactly a boundary encounter for which such a faithful lift exists --- this definition and theorem are where that possibility is constructed.

\section{Main construction theorem}

The preceding constructions now compose into the theorem that states the material-generative HCFM schema in one causal form.

\begin{theorem}[Strong material-generative HCFM theorem]
\label{thm:strong-material-generative-hcfm}
Let $\Omat$ be a stably law-indexed observation diagram in a material composition object. Assume generator regularity for every observable composite and physically realized observer, viability, and error-correction structures wherever corresponding HCFM columns are invoked. Then:
\begin{enumerate}[label=\arabic*.]
\item $\Omat$ has the canonical global material-law completion $(\Minf, \Phi^{\star})$;
\item $\Phi^{\star}$ emits the unique, natural, refinement-coherent interaction inventory $\mathrm{Emit}(C)$ on every observable composite $C$;
\item closure and quotienting derive $P_{C,H}$, and physically realized observation derives the canonical HCFM instance $S^{\circ}_{\mathrm{HCFM}}(C, H)$, with any enrichment beyond the canonical columns recorded as declared realization data;
\item every material system observable at a declared stable horizon has a
stable material-observability witness and therefore belongs to
\(\ObsMatcl\) by Theorem~\ref{thm:observability-closure}; the
variance-aware schema assignment supplies its same composition-closed canonical
HCFM realization, while the geometry, presentation, refinement, and
finite-order lift modules preserve that realization at their declared
resolutions;
\item compatible horizons have the canonical completion $P_\infty$, finite task slices have faithful coherent flat $\mathbb{R}^3$ carriers, and observable invariants transport along typed dynamic-equivariant bridges.
\end{enumerate}
\end{theorem}

\begin{proof}
Item~1 is Theorem~\ref{thm:global-material-completion}. Item~2 is
Theorem~\ref{thm:canonical-interaction-extraction}. Item~3 is
Theorems~\ref{thm:least-continuation-complete-domain},
\ref{thm:physical-functional-quotient}, and
\ref{thm:derived-observer-columns}. Item~4 is
Theorems~\ref{thm:observability-closure},
\ref{thm:hcfm-canonical-core-bridge},
\ref{thm:composition-closure},
\ref{thm:presentation-factorization}, and
\ref{thm:refinement-meta-lift-closure}. Item~5 is
Theorems~\ref{thm:horizon-completion} and \ref{thm:faithful-3d-carrier},
together with the invariant-factorization and invariant-transport results.
\end{proof}

The three modules that follow consume this theorem's outputs and develop what the completed construction makes measurable, transportable, and controllable: the dimensionality module consumes a realized registration quotient and declared metric data and returns a measured effective dimension; the disciplinary module consumes the schema's conceptual, invariant, and transport structure and returns an auditable architecture for cross-field reasoning and cumulative solving; the alignment module consumes windows, lifts, and the mass invariant and returns a conditional control picture for capability growth. Each module declares its own added premises where it uses them and closes with the ledger that types every one of its claims.

\section{Registration dimension, dimensional resolution, and the modal boundary}

Physically realized observation has produced a registration quotient; dimensionality now becomes a measurable property of it. This section defines the cost, scale, and viability data that the measurement consumes, computes the finite-scale registration dimension, and proves the conditions under which it resolves to three.

\begin{definition}[Costed registration datum]
For an observer-containing system $(C, H) \in \ObsMatcl$, a costed registration datum is
\[
D_{C,H} = \big(\widehat{A}_{C,H}, d_b, \widehat{V}_{C,H}, W_H, \tau_H\big),
\]
where $(\widehat{A}_{C,H}, d_b)$ is the metric quotient obtained by applying Definition 34 to the registration path category with a declared registration-label cost; $V_{C,H} \subseteq \widehat{A}_{C,H}$ is the viable accessed region; $W_H$ is a nonempty scale ladder: a finite set of measuring triples $(x, R, r)$ with $x \in V_{C,H}$ and $0 < r < R$; and $0 < \tau_H < \tfrac{1}{2}$ is a declared resolution tolerance. Require every $V_{C,H} \cap B_{d_b}(x, R)$ on the ladder to be totally bounded.

For $(x, R, r) \in W_H$, let $N_V(x; R, r)$ be the least number of $d_b$-balls of radius $r$, with centres in $V_{C,H}$, that cover $V_{C,H} \cap B_{d_b}(x, R)$, and set
\begin{equation}
e_H(x; R, r) = \frac{\log N_V(x; R, r)}{\log(R/r)}, \qquad I_H = \Big[\inf_{W_H} e_H,\ \sup_{W_H} e_H\Big]. \tag{S16}
\end{equation}
\end{definition}

\begin{definition}[Finite-scale effective registration dimension]
A costed registration datum is dimensionally resolved when there is a unique $d \in \mathbb{N}_0$ for which
\begin{equation}
I_H \subset (d - \tau_H, d + \tau_H). \tag{S17}
\end{equation}
In that case define $d_A(H) = d$. If no such unique integer exists, define $d_A(H) = \bot$, and call the horizon dimensionally unresolved. $d_A$ is a horizon-indexed reported coarse dimension: the number of independent directions of difference this observer resolves, at this precision, on its viable accessed region.
\end{definition}

\noindent\plainterms{} Cover a patch of the observer's tell-apart map with balls of a smaller radius and count how many you need; the way that count grows as the balls shrink is the dimension the observer actually registers. A line needs proportionally more balls (dimension 1), a sheet needs the square (dimension 2), and so on. The ladder $W_H$ is just the declared set of patch-and-ball sizes at which the count is taken, and the tolerance says how cleanly the growth exponent must sit near a whole number before a dimension is reported. If it does not sit cleanly, the honest answer $\bot$ is reported instead of a guess.

\begin{definition}[Admissible and realized registration moduli classes]
A bare admissible registration architecture is an isomorphism class of typed registration data $(A, S_A, R_A, W_A)$ satisfying the registration-side typing conditions of Definition 19, but with no selected material composite, observer, horizon, probe family, readout, label cost, viable accessed region, scale ladder, or tolerance. Let
\begin{equation}
M^{\mathrm{adm}}_{\mathrm{reg}} = \{[A]_{\cong} : A \text{ is a bare admissible registration architecture}\}. \tag{S18}
\end{equation}
The realized registration moduli class is
\begin{equation}
M^{\mathrm{real}}_{\mathrm{reg}} = \{[A]_{\cong} \in M^{\mathrm{adm}}_{\mathrm{reg}} : \exists (C, H) \in \ObsMatcl \text{ whose realized registration quotient has bare architecture } A\}. \tag{S19}
\end{equation}
``Pre-registration'' here is modal and architectural rather than chronological: $M^{\mathrm{adm}}_{\mathrm{reg}}$ records forms that a realization can select, and acquires topology, measure, or dynamics only when a future extension supplies them.
\end{definition}

\begin{proposition}[Type separation]
The moduli classes $M^{\mathrm{adm}}_{\mathrm{reg}}$ and $M^{\mathrm{real}}_{\mathrm{reg}}$ are distinct in type from both the mathematical sort $M$ in Definition 19 and the material completion $\Minf$ in Theorem 7. A bare moduli element is an isomorphism class of possible registration architecture; a mathematical-sort object is a component induced within a realized HCFM instance; and $\Minf$ is an inverse limit of material carriers. The realization map $(C, H) \mapsto [A_{C,H}]_{\cong}$ has codomain $M^{\mathrm{real}}_{\mathrm{reg}}$, and this map is the sole displayed connection among the three types.
\end{proposition}

\begin{proof}
Definition 19 places $M$ inside a realized four-sort schema and derives it through the stated bridges. Theorem 7 constructs $\Minf$ as a coherent thread of material carrier coordinates. By Definition 11.3, an element of $M^{\mathrm{adm}}_{\mathrm{reg}}$ is neither such a carrier thread nor a within-instance mathematical object: it is an isomorphism class of bare registration architecture. The sole displayed connection is the realization map to the subclass $M^{\mathrm{real}}_{\mathrm{reg}}$, which preserves rather than identifies these types.
\end{proof}

\begin{remark}[Where dimension becomes definable]
Definition 11.2 consumes a metric quotient, viable accessed region, scale ladder, and tolerance, so $d_A$ becomes definable exactly where those data exist: on declared realized costed registration data. A bare element of $M^{\mathrm{adm}}_{\mathrm{reg}}$ contains none of the selected data, so the well-typed questions there are which admissible architectures are realized and, once a realization supplies the missing data, what $d_A$ reports along its refinement thread. This placement of dimension --- as a derived, scale-dependent observable of a realized measurement rather than an antecedent background label --- is the construction's native form of a stance several quantum-gravity programs arrive at independently [12, 7, 9].
\end{remark}

\begin{proposition}[Invariance under costed registration isomorphism]
Let $F : D_{C,H} \to D_{C',H'}$ be an isometry of metric quotients that maps the viable accessed region, scale-ladder triples, and tolerance of the source to those of the target. Then $I_H = I_{H'}$, and hence $d_A(H) = d_A(H')$, including simultaneous unresolved status.
\end{proposition}

\begin{proof}
An isometry maps every eligible radius-$r$ cover of each source ball to an eligible cover of the corresponding target ball and conversely. Hence the two covering numbers, their logarithmic exponents, and the endpoint interval in (S16) agree. Definition 11.2 then gives the same reported value or unresolved status.
\end{proof}

\begin{proposition}[Refinement-indexed dimension flow]
Let $\gamma = (H_0 \preceq H_1 \preceq \cdots)$ be a refinement thread equipped at each stage with a costed registration datum. Its dimension-flow profile is the partial sequence
\begin{equation}
D_\gamma(i) = d_A(H_i) \in \mathbb{N}_0 \cup \{\bot\}. \tag{S20}
\end{equation}
Costed-registration-isomorphic refinement threads have identical profiles. The profile is a pointwise report: any monotonicity, convergence, or substrate interpretation is a further empirical or model-specific finding about a particular thread.
\end{proposition}

\begin{proof}
The profile is defined pointwise by Definition 11.2. The conclusion for equivalent threads follows at each stage from Proposition 11.6. The construction introduces neither an order relation on reported values nor a measure on threads, so further dynamical claims require further data.
\end{proof}

\begin{remark}[Relation to spectral-dimension flow]
The profile (S20) is a native registration-level analogue of scale-dependent dimension flow: it is computed from protocol-costed covering profiles of a viable operational quotient, where a diffusion spectral dimension is computed from a random walk on a geometry. The kinship is the shared stance that dimension can run with scale; each estimator retains its own definition and interpretation [7, 8, 9].
\end{remark}

\begin{proposition}[Composition in the uncoupled product regime]
Suppose the viable accessed registration datum of a composite is, over its declared scale ladder, a metric product $V_C \times V_D$ with the max product metric, product ladders, and independent covering profiles. If $d_A(H_C) = d_C$ and $d_A(H_D) = d_D$ are resolved with tolerances chosen so that the product profile is resolved, then
\begin{equation}
d_A(H_{C \otimes D}) = d_C + d_D. \tag{S21}
\end{equation}
This addition law is exactly as wide as its uncoupled premise: emitted interactions can alter the viable quotient and thereby the reported dimension, which is what the coupled-clock witness computes.
\end{proposition}

\begin{proof}
For max-metric balls on the product ladder, independent covers give matching upper and lower product-cover bounds. Their logarithms add, so the exponent interval is the sum of the component intervals. The resolution condition then selects the integer $d_C + d_D$.
\end{proof}

\begin{definition}[Separating rank, saturation, and exhaustion]
A rank-$k$ separating protocol family on $V_{C,H}$ is a family of $k$ protocol coordinates $\xi = (\xi_1, \dots, \xi_k)$ such that each coordinate is individually resolvable above the declared tolerance and, for each coordinate, there are viable pairs separated by that coordinate while the remaining coordinates are fixed to the declared resolution. Let $\mathrm{srank}(H)$ be the largest certified $k$, if one exists. The horizon is saturated at rank $k$ when $\mathrm{srank}(H) \geq k$: all $k$ directions are exercised non-degenerately. It is exhausted at rank $k$ when it has a certified rank-$k$ family and every admissible separating protocol factors through that family up to uniformly bounded resolution residue on $V_{C,H}$: no further separating direction clears the threshold.
\end{definition}

\begin{definition}[Rank--metric regularity]
A rank-$k$ separating family $\xi$ is rank--metric regular on $W_H$ when, at every $(x, R, r) \in W_H$, its resolved coordinate image contains and is contained in a $k$-fold product of chain or cycle segments whose covering numbers obey constants $0 < a_H \leq b_H < \infty$:
\begin{equation}
a_H(R/r)^k \leq N_V(x; R, r) \leq b_H(R/r)^k, \tag{S22}
\end{equation}
and the declared ladder is sufficiently separated that
\begin{equation}
\frac{\max\{|\log a_H|, |\log b_H|\}}{\inf_{W_H} \log(R/r)} < \tau_H. \tag{S23}
\end{equation}
This is the typed bridge from protocol rank to the metric-covering estimator: an additional empirical or model-specific certificate that a particular realization must supply.
\end{definition}

\begin{lemma}[Lower covering bound]
Under rank--metric regularity at rank $k$, $\inf I_H > k - \tau_H$.
\end{lemma}

\begin{proof}
The lower inequality in (S22) gives $e_H(x; R, r) \geq k + \log a_H / \log(R/r)$. Condition (S23) bounds the magnitude of the correction by $\tau_H$ uniformly over the ladder. Taking the infimum gives the claim.
\end{proof}

\begin{lemma}[Upper covering bound]
Under rank--metric regularity at rank $k$, $\sup I_H < k + \tau_H$.
\end{lemma}

\begin{proof}
The upper inequality in (S22) gives $e_H(x; R, r) \leq k + \log b_H / \log(R/r)$. Condition (S23) uniformly bounds the correction by $\tau_H$. Taking the supremum gives the claim.
\end{proof}

\begin{proposition}[Conditional operator-count agreement]
Suppose a rank--metric regular horizon has an inverse-closed registration basis with exactly one cost-separated, independently resolvable generator for each coordinate of a certified separating family, and every other registered generator is a bounded-resolution composition of these. Then the number of independent basis generators, $\mathrm{srank}(H)$, and $d_A(H)$ agree whenever the latter is resolved.
\end{proposition}

\begin{proof}
The generator conditions identify independent operational directions with the certified coordinate family, so the basis count equals $\mathrm{srank}(H)$. Rank--metric regularity bounds the exponent interval around that same rank; resolution selects it by Definition 11.2.
\end{proof}

\begin{theorem}[Estimator-calibration theorem: registration dimension at certified rank three]
Let $(C, H)$ be an observer-containing system with a declared costed registration datum. Assume:
\begin{enumerate}[label=(S\arabic*)]
\item[(S1)] Viability: $V_{C,H}$ supports persistent bounded configurations sufficient to retain the observer's declared probe and readout capability.
\item[(S2)] Self-calibration: the observer implements the strictly contractive correction structure of Theorem 22 by comparing predicted and registered continuations at the declared resolution.
\item[(S3)] Saturation: $H$ is saturated at rank 3 in the sense of Definition 11.10.
\item[(S4)] Exhaustion: $H$ is exhausted at rank 3 in the sense of Definition 11.10.
\item[(S5)] Rank--metric regularity: the certified rank-3 family is rank--metric regular on $W_H$ in the sense of Definition 11.11.
\end{enumerate}
Then $H$ is dimensionally resolved and $d_A(H) = 3$.
\end{theorem}

\begin{proof}
Premises (S1)--(S2) type the theorem's domain as viable self-calibrating observer horizons. Premise (S3) supplies the three certified independent protocol directions, and (S4) states that no additional separating direction remains outside their bounded-resolution residue. Premise (S5) supplies the metric-covering bridge at exactly this certified rank. Lemmas 11.12 and 11.13, with $k = 3$, give $I_H \subset (3 - \tau_H, 3 + \tau_H)$. Since $\tau_H < \tfrac{1}{2}$, the integer selected in Definition 11.2 is uniquely 3.
\end{proof}

\begin{proposition}[The estimator reads certified rank]
If a horizon possesses a certified separating family of rank $k$ that is rank--metric regular on its declared ladder, then $I_H \subset (k - \tau_H, k + \tau_H)$ and $d_A(H) = k$.
\end{proposition}

\begin{proof}
Apply Lemmas 11.12 and 11.13 at rank $k$. The tolerance windows around distinct integers are disjoint because $\tau_H < \tfrac{1}{2}$.
\end{proof}

\begin{corollary}[Rank-three premise sharpness]
Under rank--metric regularity of the maximal certified separating family, saturation at rank three rules out every lower reported rank and exhaustion at rank three rules out every higher certified rank. Thus saturation and exhaustion fix the rank that Proposition 11.16 converts to the reported dimension.
\end{corollary}

\begin{remark}[No bridge, no verdict]
Saturation and exhaustion without rank--metric regularity need not determine $d_A$. For example, a viable three-coordinate product with a declared ladder spanning fine scales that resolve all three factors and coarser scales that saturate one factor has covering exponents near both three and two; its interval $I_H$ crosses no unique integer window and the honest output is $\bot$. The rank--metric bridge is therefore the estimator premise, while saturation and exhaustion select the certified rank to which it is applied.
\end{remark}

\begin{remark}[Logical structure of the calibration theorem]
Theorem 11.15 is a soundness-and-calibration result for the registration-dimension instrument. Rank--metric regularity supplies the covering estimate; saturation and exhaustion fix its certified rank; viability and self-calibration delimit the intended realization class and are consumed by the genericity conjecture and its empirical discharge paths.
\end{remark}

\noindent\plainterms{} Read the theorem's premises as the two conditions under which an observer's registered dimension is forced to match a count of three: saturation says all three independent directions of difference are genuinely exercised (none is degenerate), and exhaustion says no fourth direction clears the observer's threshold --- ``three'' is then the rank of the observer's fitness-relevant registration map, and the rank--metric certificate converts that rank into the ball-counting exponent, pinning the measured value at exactly 3. The result is a statement about a certified observer, delivered with its certificates listed --- which is what makes it testable rather than rhetorical.

\begin{remark}[Filter honesty and the role of classical dimensionality arguments]
Theorem 11.15 derives a reported registration dimension of three from its declared saturation, exhaustion, and rank--metric premises. Which realizations satisfy those premises is the separate question the genericity conjecture below addresses. Classical analyses of orbital and atomic stability and of predictive well-posedness [2, 3, 4] are contextual constraints on possible premise certification --- they motivate tests of which large-dimensional regimes can support familiar persistent observers under particular law families --- and the theorem's viability and self-calibration premises are where that classical selection pressure enters this construction in typed form.
\end{remark}

\begin{conjecture}[Restricted genericity of infrared registration dimension]
Let a declared realization supply a reachable subclass $M^{\mathrm{reach}}_{\mathrm{reg}} \subseteq M^{\mathrm{real}}_{\mathrm{reg}}$, a probability measure $\mu$ on specified realization/refinement trajectories, and a dynamics $G$ preserving viable, maximally separating, self-calibrating costed registration data. The conditional prediction is that $\mu$-typical trajectories contain an infrared horizon satisfying the premises of Theorem 11.15. Equivalently, the predicted generic infrared report on that declared restricted domain is $d_A = 3$.
\end{conjecture}

\begin{remark}[Globally coherent genericity route]
Theorem 11.15 supplies the consequence of the declared realization conditions: whenever the typical trajectories specified in Conjecture 11.21 reach a certified infrared horizon, their report is $d_A = 3$. The genericity claim is therefore a globally coherent conditional prediction on every realization class that supplies the stated reachable domain, measure, dynamics, and certification route. The discharge paths below determine which material realization classes instantiate that prediction; they do not remove its conditional content from the present construction.
\end{remark}

\begin{remark}[Preregistered discharge path DP1: causal-order bridge]
A causal-order route supplies a typed bridge from a realized registration datum to a causal-order estimator and proves the conversion condition $d_A = d_{MM} - 1$ in the declared Lorentzian-registration regime, where $d_{MM}$ is the Myrheim--Meyer order-dimension estimate. Outcomes are typed in advance: (O1) a certified bridge with generic $d_{MM} = 4$ realizes the restricted genericity prediction; (O2) a certified bridge with a non-four-dimensional generic estimate, or a failed bridge, identifies a realization class that does not instantiate it; (O3) absence of a bridge or stable estimate leaves that realization class unclassified. Causal-set dimension estimators motivate this route and supply its target quantity [10, 11, 12].
\end{remark}

\begin{remark}[Preregistered discharge path DP2: ultraviolet departure profile]
A refinement route measures the profile $D_\gamma$ on a declared family of increasingly fine costed registration data. A reproducible departure from the infrared value at shorter registered scales, together with an independently certified return to the conditions of Theorem 11.15 in the infrared, realizes a nontrivial regime structure for the genericity prediction; a profile inconsistent with the declared outcome criterion identifies a realization class outside that regime; failure to obtain a resolved profile leaves that class unclassified. The posture is that of dimensional-flow programs, executed at the registration level [7, 8, 9].
\end{remark}

\begin{remark}[Module ledger and scope]
Definitions 11.1--11.3, Propositions 11.4--11.14, and Theorem 11.15 are complete within their stated premises; Remark 11.22 records the positive conditional genericity prediction and its two realization paths. A measure, dynamics, or topology on $M^{\mathrm{adm}}_{\mathrm{reg}}$, a causal-order/registration bridge, a dimension assignment to $\Minf$, and a carrier--registration identification are typed as realization data for particular empirical domains, not as dependencies of the preceding theorems. Their construction determines which domains instantiate the globally coherent prediction.
\end{remark}

\section{Disciplinary realizations, perspective-explicit transport, and reusable semantic search}

The completed construction contains conceptual, registration, mathematical, and fitness structures with typed bridges. This module uses them to construct what a field of study is, what must travel with a piece of reasoning for it to stay valid outside the field that produced it, and how certified solutions accumulate into a compounding problem-solving capacity.

The founding statement records that the same modeling approach has been used to model physics, cognition, consciousness, and other systems. This module supplies the interface that makes such applications auditable and their results transportable: a discipline is specified as a locally bounded realization whose objects and reasoning processes are meaningful relative to an explicit perspective, evaluator, admissibility relation, semantic equivalence, and certificate discipline. Any embedding of the disciplines into one completed human conceptual space, if later constructed, must preserve exactly this interface --- which is why the module's results are stated for discipline-local interfaces and remain meaningful whether such a completion is eventually built, partially realized, or left open.

\subsection{The disciplinary object and perspective-explicit reasoning}

\begin{definition}[Disciplinary realization]
A disciplinary realization is a typed tuple
\begin{equation}
D = \big(\Omega_D, X_D, \mathrm{Op}_D, \mathrm{Adm}_D, \mathrm{Eval}_D, \equiv_D, \mathrm{Cert}_D, \mathrm{Persp}_D\big), \tag{S29}
\end{equation}
where:
\begin{itemize}
\item $\Omega_D$ is the type of domain objects and $X_D$ is the associated state or representation class;
\item $\mathrm{Op}_D$ is a typed path category of admissible primitive and composite reasoning operations on $X_D$;
\item $\mathrm{Adm}_D \subseteq \mathrm{Path}_D$ is the predicate selecting paths admissible under the discipline's declared constraints;
\item $\mathrm{Eval}_D$ assigns declared outcome, input, fitness, error, or other evaluation values to admissible paths in a specified codomain;
\item $\equiv_D$ is a semantic equivalence on admissible paths or their endpoint-labelled presentations;
\item $\mathrm{Cert}_D(p)$ is the certificate type required to establish that a path $p$ is admissible and that its declared evaluation is valid within scope; and
\item $\mathrm{Persp}_D = (O_D, H_D, W_D, Q_D)$ is the perspective datum: observer or evaluator role, horizon or resolution, fitness/viability context, and question/readout class under which the preceding components are interpreted.
\end{itemize}
A realization of $D$ in an HCFM schema is a declared typed interface from its objects, operations, evaluator, and certificates to the relevant $A, C, M, W$ columns of a specified $S_{\mathrm{HCFM}}(C, H)$. A realization in a supplied global conceptual interface $K$, if one exists, is a separate map with its own certificate.
\end{definition}

\begin{definition}[Perspective-implicit and perspective-explicit process]
A first-order process in $D$ is an admissible path $p : x \rightsquigarrow y$ in $\mathrm{Op}_D$ acting on domain objects. It is perspective-implicit when the data in $\mathrm{Persp}_D$, $\mathrm{Adm}_D$, $\mathrm{Eval}_D$, $\equiv_D$, or $\mathrm{Cert}_D$ are required for its validity but are not carried in the stated process interface.

A second-order disciplinary process is a tuple
\begin{equation}
p_e = (p; \mathrm{Persp}_D, \mathrm{Adm}_D, \mathrm{Eval}_D, \equiv_D, \mathrm{Cert}_D), \tag{S30}
\end{equation}
whose domain and codomain make those conditions available for inspection: a perspective-explicit interface through which the conditions of valid inference can themselves be compared, transformed, or transported.
\end{definition}

\begin{proposition}[Native validity does not imply transport validity]
An admissible first-order process $p$ in $D$ has no well-typed cross-disciplinary validity claim merely from its existence as a path in $\mathrm{Op}_D$. A transport claim to another disciplinary realization $D'$ requires a map that specifies the image of the process and the preservation or translation of its perspective, admissibility, evaluator, semantic-equivalence, and certificate conditions.
\end{proposition}

\begin{proof}
The predicate $\mathrm{Adm}_D$, evaluation $\mathrm{Eval}_D$, relation $\equiv_D$, and certificate type $\mathrm{Cert}_D$ have source discipline $D$ in their type. A path in $\mathrm{Op}_D$ alone contains no target object in $\mathrm{Op}_{D'}$, no map between the predicates or evaluators, and no certificate conversion. Thus the expression ``$p$ remains valid in $D'$'' is not typed until those maps are declared.
\end{proof}

\begin{remark}[First-order rigor and the transport-fragile failure mode]
A first-order process can be fully rigorous within its native discipline. The term transport-fragile identifies the specific additional failure mode that appears when a process whose perspective conditions were locally sufficient is reused across a different evaluator, horizon, object type, or fitness relation without exposing the bridge that would preserve those conditions --- the proposition locates the failure at the missing bridge, not in the first-order work.
\end{remark}

\subsection{Perspective-explicit bridges and disciplinary invariant transport}

\begin{definition}[Perspective-explicit disciplinary bridge]
Let $D$ and $D'$ be disciplinary realizations. A perspective-explicit bridge
\begin{equation}
F_{D \to D'} = (F_X, F_{\mathrm{Op}}, F_{\mathrm{Persp}}, F_{\mathrm{Eval}}, F_{\equiv}, F_{\mathrm{Cert}}) \tag{S31}
\end{equation}
is a family of typed maps such that:
\begin{enumerate}[label=(B\arabic*)]
\item[(B1)] $F_X$ maps source objects and $F_{\mathrm{Op}}$ is a functor mapping admissible source paths to target paths while preserving composition and identities;
\item[(B2)] $F_{\mathrm{Persp}}$ maps or explicitly restricts the observer, horizon, question, and viability data;
\item[(B3)] $p \in \mathrm{Adm}_D$ implies $F_{\mathrm{Op}}(p) \in \mathrm{Adm}_{D'}$, or a declared target residue records the exact condition under which admissibility fails;
\item[(B4)] $F_{\mathrm{Eval}}$ preserves the declared evaluator relation on the bridge's stated scope, or records a calibrated translation with its residual;
\item[(B5)] $p \equiv_D q$ implies $F_{\mathrm{Op}}(p) \equiv_{D'} F_{\mathrm{Op}}(q)$; and
\item[(B6)] $F_{\mathrm{Cert}}$ maps a valid source certificate to a target certificate or to an explicitly typed obligation still requiring discharge.
\end{enumerate}
A bridge is certified when all six clauses hold on its declared domain.
\end{definition}

\begin{proposition}[Disciplinary specialization of invariant transport]
Let $F_{D \to D'}$ be a certified perspective-explicit bridge whose underlying operation map is a typed dynamic-equivariant bridge between the relevant HCFM schema realizations. Let $I_{D'}$ be an observable property that factors through $\equiv_{D'}$ and whose evaluator scope is preserved by $F_{\mathrm{Eval}}$. Then $I_{D'} \circ F_{\mathrm{Op}}$ factors through $\equiv_D$ and is an observable invariant of the source realization. If the bridge is an equivalence with inverse certificate translation, the invariant transport is bijective.
\end{proposition}

\begin{proof}
By clause (B5), semantically equivalent source paths have semantically equivalent target images. Since $I_{D'}$ factors through target equivalence, the composite is constant on source equivalence classes and therefore factors through $\equiv_D$. Clauses (B2)--(B4) ensure that the target evaluator and scope invoked by $I_{D'}$ remain applicable on the bridge's declared domain; clause (B6) preserves the evidence condition. Dynamic equivariance supplies the underlying HCFM invariant-transport condition of Proposition 30. An equivalence provides the inverse argument.
\end{proof}

\begin{obligation}[Bridge audit]
A proposed cross-disciplinary reuse must record: source and target objects; source and target perspectives; the operation functor; admissibility and evaluator residuals; semantic-equivalence preservation; certificate translation; and the domain on which each is checked. A surface analogy, common vocabulary, or shared drawing is not a substitute for this record.
\end{obligation}

\noindent\plainterms{} A discipline is a workshop; second-order work ships each tool with its label --- what it is for, when it is safe, how to check the result, from whose vantage point it was calibrated. The bridge (B1)--(B6) is the checklist for moving a tool to another workshop with its label intact, and the transport proposition is the payoff: a property established in the target workshop pulls back to a genuine property in the source workshop whenever the bridge is certified. The audit obligation says what everyone quietly knows: a shared word or a suggestive diagram is not a bridge.

\subsection{The observer-inclusion blind spot, recursive global coherence, and stop conditions}

There is an observer-inclusion blind spot in the general assessment of systems when a hypothesis is judged relative to a fixed observer, horizon, admissibility relation, evaluator, semantic-equivalence relation, certificate discipline, or transport rule while one or more of those assessment conditions is not represented in the system model. The observed/observing-system distinction is a classical reason to make the observer explicit [14]. The claim here is narrower and operational: in a novel system or novel problem-domain, a locally adequate one-shot assessment can be insufficient when the assessment conditions themselves are among the variables whose joint satisfiability determines whether the hypothesis is coherent.

A local-prior assessment is appropriate when familiar evidence, stable regularities, and a strongly constrained system model fix the relevant assessment conditions. The blind spot is instead a risk class for systems with sparse priors, underdetermined model variables, or uncertified cross-system transport. The literature supplies formal analogues rather than a substitute for the definitions below. Causal identifiability distinguishes targets determined by observational structure from targets on which observationally compatible models can disagree [15]; partial-identification analysis retains an identified set when the available data and assumptions do not select one value [16]. Infinite belief hierarchies show that perspective structure can have no a priori finite depth [17]. These results motivate explicit assessment conditions; the HCFM blind-spot diagnosis and global-coherence protocol are defined by the objects and certificates introduced here.

\begin{definition}[HCFM assessment datum and local-prior assessment]
Let $S$ be a system represented by a declared HCFM system model $M_0(S)$, and let $h$ be a hypothesis about $S$. An HCFM assessment datum is a tuple
\begin{equation}
A_0(S, h) = \big(M_0(S), h, \Xi_0, B, \eta, \kappa, G\big), \tag{S32a}
\end{equation}
where
\[
\Xi_0 = (O, H, W, Q, \mathrm{Adm}, \mathrm{Eval}, \equiv, \mathrm{Cert})
\]
is the declared assessment perspective, $B$ is the coherence budget, $\eta$ is a declared assessment noise floor, $\kappa$ is an observability ceiling, and $G$ is a declared hypothesis-and-model extension generator. The model $M_0(S)$ is continuous on its declared system variables when its material or operational evolution is given by continuous maps on the relevant carrier and the permitted parameter or horizon refinements are continuous in the declared topology.

A local-prior assessment fixes $(M_0, \Xi_0)$ and evaluates $h$ against the evidence, priors, and certificates available at that initial horizon. A recursive global-coherence correction is a certified transition
\begin{equation}
(M_n, \Xi_n) \rightsquigarrow (M_{n+1}, \Xi_{n+1}) \tag{S32b}
\end{equation}
which adds, restricts, or recalibrates a model component, observer condition, horizon, admissibility condition, evaluator, semantic-equivalence relation, certificate rule, or transport bridge. Continuity makes controlled perturbation and refinement available in one declared model class; the additional generator $G$, rather than continuity alone, supplies the potentially unbounded family of falsifiable hypotheses and model extensions.
\end{definition}

\subsubsection*{Proof of the recursive non-factorization theorem}

Fix an HCFM assessment datum $A_0(S,h)$ with local-prior perspective
$(M_0,\Xi_0)$, and let $\mathcal{X}$ be a declared class of candidate
completions of that datum. Let
\[
\Lambda_L : \mathcal{X} \longrightarrow \mathcal{R}_L
\]
be complete for the assessments admitted while $(M_0,\Xi_0)$ remains fixed:
every locally admitted observation, prior, certificate, and result of a
recursively generated local-prior obligation is determined by
$\Lambda_L(x)$. Let
\[
\Omega : \mathcal{X} \longrightarrow \mathcal{Q}
\]
be a declared target of global assessment.

An $L$-confined recursive assessment is any sequence of assessment states
$\{z_n\}_{n\geq 0}$ for which there exist maps $a_0,a_1,\ldots$ satisfying
\[
z_0(x)=a_0\!\left(\Lambda_L(x)\right)
\]
and
\[
z_{n+1}(x)
=
a_{n+1}\!\left(
\Lambda_L(x),
z_0(x),\ldots,z_n(x)
\right).
\]
Thus the recursion may generate arbitrarily many child obligations, but it
remains confined to the distinctions retained by $\Lambda_L$.

\noindent\textbf{Claim.}
There exists a map
\[
\overline{\Omega}_L :
\operatorname{im}(\Lambda_L) \longrightarrow \mathcal{Q}
\]
such that
\[
\Omega=\overline{\Omega}_L\circ\Lambda_L
\]
if and only if $\Omega$ is constant on every fiber of $\Lambda_L$.
If there exist $x,y\in\mathcal{X}$ such that
\[
\Lambda_L(x)=\Lambda_L(y)
\qquad\text{and}\qquad
\Omega(x)\neq\Omega(y),
\]
then every $L$-confined recursive assessment produces the same assessment
history on $x$ and $y$. Therefore no such assessment can determine
$\Omega$ correctly on both candidates. Any record
\[
\Lambda_G : \mathcal{X} \longrightarrow \mathcal{R}_G
\]
through which $\Omega$ is correctly determined on both candidates must
satisfy
\[
\Lambda_G(x)\neq\Lambda_G(y).
\]

\begin{proof}
Suppose first that
\[
\Omega=\overline{\Omega}_L\circ\Lambda_L.
\]
If $\Lambda_L(x)=\Lambda_L(y)$, then
\[
\Omega(x)
=
\overline{\Omega}_L\!\left(\Lambda_L(x)\right)
=
\overline{\Omega}_L\!\left(\Lambda_L(y)\right)
=
\Omega(y),
\]
so $\Omega$ is constant on every fiber of $\Lambda_L$.

Conversely, suppose that $\Omega$ is constant on every fiber of
$\Lambda_L$. Define
\[
\overline{\Omega}_L(r):=\Omega(x)
\]
for any $x\in\mathcal{X}$ satisfying $\Lambda_L(x)=r$. This is well-defined:
if both $x$ and $y$ satisfy $\Lambda_L(x)=\Lambda_L(y)=r$, fiberwise
constancy gives $\Omega(x)=\Omega(y)$. By construction,
\[
\Omega=\overline{\Omega}_L\circ\Lambda_L.
\]
Uniqueness on $\operatorname{im}(\Lambda_L)$ follows because every
$r\in\operatorname{im}(\Lambda_L)$ has at least one preimage.

Now suppose that
\[
\Lambda_L(x)=\Lambda_L(y).
\]
The defining equation for $z_0$ gives $z_0(x)=z_0(y)$. If
$z_k(x)=z_k(y)$ for every $k\leq n$, then the defining equation for
$z_{n+1}$ gives
\[
z_{n+1}(x)
=
a_{n+1}\!\left(
\Lambda_L(x),z_0(x),\ldots,z_n(x)
\right)
=
a_{n+1}\!\left(
\Lambda_L(y),z_0(y),\ldots,z_n(y)
\right)
=
z_{n+1}(y).
\]
Induction therefore gives identical recursive assessment histories for $x$
and $y$. If $\Omega(x)\neq\Omega(y)$, no verdict extracted from that common
history can equal $\Omega$ on both candidates.

Finally, suppose that a record $\Lambda_G$ determines $\Omega$ on both
$x$ and $y$. If $\Lambda_G(x)=\Lambda_G(y)$, the same factorization
argument would force $\Omega(x)=\Omega(y)$, contrary to hypothesis.
Therefore
\[
\Lambda_G(x)\neq\Lambda_G(y).
\]
The resolving assessment must retain or create a distinction that the
fixed local record does not retain. In the HCFM setting, that requires a
certified extension or recalibration of representation, observation,
horizon, model class, or bridge structure.
\end{proof}

\begin{definition}[Observer-inclusion blind spot]
For an assessment datum $A_0(S, h)$, an observer-inclusion blind spot is present when a component of $\Xi_0$ is held fixed outside $M_0(S)$, while a certified recursive correction establishes that the component changes the admissible model class, target invariant, evaluator scope, bridge validity, or certificate condition required to assess $h$. The blind spot is systemic when that omitted component concerns a condition shared by a class of systems or by a transport relation among them, rather than an idiosyncratic measurement error. The blind spot is realized only when its correction residual is measurable on a declared common scale; it is not inferred merely from disagreement with empirical evidence, consensus, or a prior assessment.
\end{definition}

\begin{definition}[Three-valued global-coherence assessment]
For an assessment datum $A_0(S, h)$, the HCFM global-coherence assessment is
\begin{equation}
\mathrm{GC}_B(S, h) \in \{\mathrm{C}, \mathrm{I}, \mathrm{UB}\}, \tag{S32c}
\end{equation}
where $\mathrm{C}$, $\mathrm{I}$, and $\mathrm{UB}$ denote, respectively, coherent, incoherent, and undetermined within coherence budget.

It returns $\mathrm{C}$ when the declared HCFM system model, hypothesis, perspective conditions, and certified extensions admit a coherence witness: a jointly satisfiable completion, invariant factorization, or certified bridge of the type required by $h$. It returns $\mathrm{I}$ when a valid counter-witness falsifies a required premise, invariant condition, bridge clause, refinement condition, or empirical realization condition. It returns $\mathrm{UB}$ when the witness or counter-witness requires a model extension, observational resolution, certified bridge, or refinement depth outside $B$, below $\eta$, or beyond $\kappa$.
\end{definition}

\begin{definition}[Internalized constraint closure and external certification]
An assessor may execute a global-coherence assessment before the full formal model tower is externally serialized, provided the assessor can internally instantiate and recursively apply the tower's declared admissibility, consequence, invariant, bridge, falsification, and stopping constraints to $h$. An internal verdict is a verdict reached by that constraint execution with the same declared falsifiers, noise floor, and observability ceiling as Definition 12.10. An externally certified verdict additionally carries system models, proof objects, certificates, and traces sufficient for independent audit, falsification, transmission, and reuse. External formalization preserves and exposes the constraint algebra; it neither creates the coherence relation nor is it a prerequisite for the first execution of the assessment.
\end{definition}

\begin{remark}[Recursive assessment of an underdetermined hypothesis]
When an internally instantiated or externally represented $M_n(S)$ underdetermines $h$, the protocol does not infer either binary label from that underdetermination. A reliable escalation from $\mathrm{UB}$ first represents the target of $h$, where possible, as an observable invariant or a typed evaluator value in $M_n(S)$. It then permits projection to another system model only through a certified perspective-explicit bridge in the sense of Definition 12.5; the transported target is reassessed in the target model, and the process may iterate. If the required invariant, bridge, resolution, or certificate is unavailable, the assessment remains $\mathrm{UB}$. Thus recursive reassessment is an auditable tower of HCFM system models rather than an unrestricted demand for eventual agreement.
\end{remark}

\begin{definition}[Certified correction residual and assessment discrepancy]
For a recursive correction (S32b), let $r_n(S, h) \geq 0$ be a certified residual in a declared common scale: for example, a perspective, admissibility, evaluator, semantic-equivalence, certificate, dynamic-commutation, or invariant-preservation residual. The cumulative assessment discrepancy after $N$ corrections is
\begin{equation}
d_N(S, h) := \sum_{n=1}^{N} r_n(S, h). \tag{S32d}
\end{equation}
The discrepancy measures correction of the system model and assessment interface, not distance between the three verdict labels.
\end{definition}

\begin{proposition}[Sufficient conditions for an unbounded observer-inclusion blind spot]
Let $A_0(S, h)$ have an observer-inclusion blind spot in the sense of Definition 12.9 and admit an unbounded sequence of certified recursive global-coherence corrections. Suppose that there exist $\varepsilon > \eta$ and infinitely many correction indices $n$ for which $r_n(S, h) \geq \varepsilon$. Then $d_N(S, h)$ is unbounded along the correction sequence. In particular, if every correction has $r_n \geq \varepsilon$, then $d_N \geq N\varepsilon$. If, after some stage, $r_{n+1} \geq \lambda r_n$ for $\lambda > 1$, then the discrepancy has an exponential lower growth bound.
\end{proposition}

\begin{proof}
For any $K > 0$, choose more than $K/\varepsilon$ indices with residual at least $\varepsilon$. The corresponding partial sum exceeds $K$, so $d_N$ is unbounded. Under the uniform condition, the displayed linear bound follows directly. Under the expansive recurrence, induction gives $r_{n+k} \geq \lambda^k r_n$; summing the geometric lower bound proves exponential growth. Since each $r_n$ is, by hypothesis, the measured residue of a correction required by the omitted assessment condition, the unbounded discrepancy realizes an unbounded observer-inclusion blind spot rather than mere disagreement among verdict labels.
\end{proof}

\begin{remark}[When the blind-spot condition is substantive]
Proposition 12.14 is algebraic. Its application to a system assessment requires additional typed facts: (i) indefinitely extensible observer, horizon, bridge, or system-model structure; (ii) local indistinguishability paired with global separation of the target, for example $M \equiv_{\Xi_0} M'$ while $\mathrm{dist}(q_M, q_{M'})$ exceeds every prescribed bound; (iii) certified non-cancelling residuals; and (iv) residuals remaining above the noise floor and below the observability ceiling. Identifiability and partial-identification theory make the local-indistinguishability issue explicit [15, 16]; infinite belief hierarchies support the possibility of unbounded perspective extension [17]; the present construction adds the typed assessment and bridge residues. Without these conditions, there is no claim that disagreement with a local prior grows, let alone that a blind spot is unbounded.
\end{remark}

\begin{definition}[Certified assessment constraint chain]
Let $\Gamma_0 \subseteq \Gamma_1 \subseteq \cdots$ be the increasing chain of all soundly certified constraints accumulated by a fair recursive assessment, and set
\begin{equation}
\Gamma_\infty := \bigcup_{n \geq 0} \Gamma_n. \tag{S32e}
\end{equation}
A hypothesis is finitely separated on the incoherent side when every incoherent instance has a finite certified counter-witness in some $\Gamma_n$. A declared HCFM model class has the coherence-completion property when every finitely satisfiable constraint family in the declared class has a certified jointly satisfiable completion. A monotone assessment update is one whose surviving-candidate or constraint operator is monotone on a declared complete lattice or otherwise has a declared convergent fixed-point construction.
\end{definition}

\begin{theorem}[Conditional convergence of recursive global-coherence assessment]
Let $A_0(S, h)$ have an unbounded coherence budget and an unbounded fair assessment schedule over its declared system-model and hypothesis class. Suppose:
\begin{enumerate}[label=(A\arabic*)]
\item[(A1)] every relevant model extension, bridge obligation, and valid falsifier in the declared class is eventually considered;
\item[(A2)] all accepted certificates and counter-witnesses are sound on their declared scope;
\item[(A3)] every incoherent instance of $h$ is finitely separated by a valid counter-witness;
\item[(A4)] the declared HCFM system-model class has the coherence-completion property for $\Gamma_\infty$;
\item[(A5)] the assessment update has a declared monotone convergent fixed-point construction; and
\item[(A6)] the resulting witness or counter-witness is resolved above $\eta$ and within $\kappa$.
\end{enumerate}
Then recursive global-coherence assessment reaches a binary limit verdict: it returns $\mathrm{I}$ once a finite valid counter-witness is exposed, or $\mathrm{C}$ at the certified completion limit when $\Gamma_\infty$ is jointly satisfiable. If any displayed condition is unavailable, or the final residual lies at the noise or observability boundary, the protocol returns $\mathrm{UB}$ rather than forcing a binary conclusion.
\end{theorem}

\begin{proof}
If $h$ is incoherent, (A3) supplies a finite counter-witness and (A1) eventually exposes it; (A2) makes the resulting $\mathrm{I}$ verdict valid. If $h$ remains coherent, all finite certified constraint sets are jointly satisfiable. By (A4), $\Gamma_\infty$ admits a certified completion; (A5) supplies the convergent assessment construction, and (A2) makes its witness valid. Condition (A6) ensures that the verdict is operationally resolvable. Failure of any premise removes the inferential route to a binary verdict, so Definition 12.10 assigns $\mathrm{UB}$.
\end{proof}

\begin{remark}[Literature boundary for the convergence theorem]
Theorem 12.17 draws on recognizable structural ingredients. Gold's identification-in-the-limit result shows why hypothesis-class and evidence restrictions matter [18]. Compactness supplies a classical local-to-global satisfiability form for logical backgrounds carrying that property [19]. Tarski's fixed-point theorem and abstract interpretation provide standard monotone fixed-point forms [1, 20]; well-structured transition systems show how infinite-state classes can remain decidable under adequate order [21]. These results locate the formal roles of Conditions (A1)--(A6); HCFM specifies the observer, bridge, certificate, and realization conditions for the present assessment setting.
\end{remark}

\noindent\plainterms{} The blind spot is not that a familiar local assessment is irrational. It is that a novel system can require assessment of the assessment conditions themselves. A recursive HCFM assessment therefore carries the workshop --- observer, horizon, model, bridge, rules, and certificates --- into the field of observation. If certified leftovers keep arriving above the noise floor, the correction can grow without bound. If the declared class also has fair tests, finite refuters, a valid completion rule, and a convergent update, the process reaches a binary limit verdict. Otherwise, the protocol reports undetermined within coherence budget rather than pretending that repetition alone decides the system.

\subsection{Semantic representations, certified libraries, and semantic backpropagation}

\begin{definition}[Semantic representation and canonical library]
For a disciplinary realization $D$, its semantic representation class is
\begin{equation}
\mathrm{Sem}(D) = \mathrm{Adm}_D / \equiv_D. \tag{S32}
\end{equation}
A canonical semantic library is a persistent tuple
\begin{equation}
L_D = (\mathrm{Node}_D, \mathrm{Edge}_D, \mathrm{can}_D, \mathrm{cert}_D, \beta_D), \tag{S33}
\end{equation}
where each node is one semantic equivalence class $[p]_{\equiv_D}$, $\mathrm{can}_D$ maps any retained presentation to its unique class representative, $\mathrm{Edge}_D$ records typed admissible composition or dependency relations among classes, $\mathrm{cert}_D$ attaches the current certificate record, and $\beta_D$ is a declared operator-count or search-cost bound. The library is semantically deduplicated when no two distinct nodes represent equivalent paths under $\equiv_D$.
\end{definition}

\begin{definition}[Invariant evaluator and semantic backpropagation]
Suppose a declared path evaluator has a nonzero input quantity $I_D(p) > 0$, an outcome quantity $O_D(p)$, and an efficiency value
\begin{equation}
\eta_D(p) = \frac{O_D(p)}{I_D(p)}. \tag{S34}
\end{equation}
The evaluator is semantic-invariant when $p \equiv_D q$ implies $\eta_D(p) = \eta_D(q)$; then it descends to $\bar{\eta}_D : \mathrm{Sem}(D) \to \mathbb{R}$. More generally, any quantity used in the search is permitted only when it factors through the declared semantic quotient or carries a typed residual.

For a task $\tau = (s, t, \theta, \beta)$, where $s$ and $t$ are typed endpoint conditions, $\theta$ is a declared improvement threshold, and $\beta \leq \beta_D$, define the bounded candidate class
\begin{equation}
P_{\tau,\beta}(L_D) = \{p : s \rightsquigarrow t : p \text{ is a certified admissible composition of library classes and } \mathrm{cost}(p) \leq \beta\}. \tag{S35}
\end{equation}
Semantic backpropagation is the reverse-dependency search procedure that begins from the target condition $t$ and its required evaluation threshold, traverses the incoming certified dependency and composition edges of $L_D$, and returns the classes
\begin{equation}
\mathrm{SBP}_{D,\beta}(\tau) = \arg\max_{p \in P_{\tau,\beta}(L_D)} \bar{\eta}_D([p]_{\equiv_D}) \tag{S36}
\end{equation}
subject to the declared certificate and admissibility checks. An empty candidate set is a typed failure result.
\end{definition}

\begin{proposition}[Bounded certified retention monotonicity]
Suppose $L_D$ is finite within each declared budget and is updated only by: (i) canonical deduplication; (ii) insertion of a certified admissible class; and (iii) replacement of an existing class record only by a certificate-preserving representative with weakly greater $\bar{\eta}_D$. Then for every fixed task $\tau$ and budget $\beta$, the best certified retained value
\begin{equation}
\eta^{*}_{D,\tau,\beta} = \max\{\bar{\eta}_D([p]) : p \in P_{\tau,\beta}(L_D)\} \tag{S37}
\end{equation}
is nondecreasing under accepted updates whenever the candidate set is nonempty.
\end{proposition}

\begin{proof}
Canonical deduplication preserves semantic classes and therefore preserves $\bar{\eta}_D$. An insertion enlarges the candidate set or leaves it unchanged. A permitted replacement retains the class and does not lower its stored evaluator value. Thus the maximum over the finite candidate set cannot decrease.
\end{proof}

\begin{remark}[What the retention proposition establishes]
Proposition 12.21 establishes that under the declared bounded, certified retention rule, the shelf of solved routes only improves: accepted updates never lose a best certified answer. Global optimality, universal semantic invariance, search cost, and empirical performance gain are each separate constructions, and the definition's value is that every one of them now has a named location to be built and tested: quotient construction, evaluator factoring, candidate generation, certificate validation, or budget control.
\end{remark}

\subsection{Compositional problem solving and combinatorial candidate capacity}

\begin{definition}[Problem, solution, and reusable module family]
A well-formed problem relative to $D$ is a tuple $\pi = (s, t, \Gamma)$ of typed source condition $s$, target condition $t$, and declared constraints $\Gamma$, for which no currently retained certified admissible path from $s$ to $t$ is known within the stated budget. A solution is a certified admissible composition $p : s \rightsquigarrow t$ satisfying $\Gamma$.

A reusable module family of size $n$ is a set $M_n = \{m_1, \dots, m_n\}$ of certified semantic classes with a common typed interface and identity module $\mathrm{id}$, such that every declared permitted composite of members is an admissible library path. The family is nondegenerate when distinct declared selections of modules produce semantically distinct certified compositions.
\end{definition}

\begin{definition}[Conditions for compositional reuse]
A persistent library has the combinatorial-reuse conditions when:
\begin{enumerate}[label=(C\arabic*)]
\item[(C1)] Composable objects: solutions are typed objects with declared composition and interface conditions, not merely stored text or traces.
\item[(C2)] Common coordinates: problem definitions, solution objects, and their endpoints are represented in a shared semantic coordinate system or through certified coordinate bridges.
\item[(C3)] Perspective-explicit transport: every cross-disciplinary reuse is mediated by a certified bridge in the sense of Definition 12.5.
\item[(C4)] Semantic deduplication and bounded search: meaning-equivalent entries collapse through $\mathrm{can}_D$, and every search reports its operator bound $\beta$.
\item[(C5)] Nondegenerate reusable interfaces: the library contains a growing module family whose certified compositions remain admissible and semantically distinct on a declared problem family.
\item[(C6)] Certificate-preserving retrieval, verification, and task realization: an indexed retrieval process returns a candidate together with a certificate whose validity can be checked in the target realization; declared task descriptors are coherently associated with certified solution classes; and the amortized work required to add a module, maintain the index, and maintain the compositional certificate scheme is recorded explicitly.
\end{enumerate}
Each condition earns its place by what its absence costs. Without (C1) the library is a lookup table: stored answers, no composition. Without (C2) it is disconnected additive fragments. Without (C3), false compositions corrupt reach. Without (C4), meaning-equivalent entries multiply and storage grows combinatorially for linear problem-solving gain --- the worst case. Conditions (C1)--(C4) state the representational and transport architecture; (C5)--(C6) convert stored composability into certified task-indexed capacity and make the rate theorem below well typed.
\end{definition}

\begin{theorem}[Conditional combinatorial certified-candidate capacity]
Let $L_D$ satisfy Definition 12.24. Suppose that for each $n$ it contains a nondegenerate reusable module family $M_n = \{m_1, \dots, m_n\}$ with a common interface such that, for every subset $I \subseteq [n]$, the ordered composition
\begin{equation}
m_I = \mathop{\circ}_{i \in I} m_i \tag{S38}
\end{equation}
(with $m_\emptyset = \mathrm{id}$) is certified admissible and $I \neq J$ implies $[m_I]_{\equiv_D} \neq [m_J]_{\equiv_D}$. Then the library contains at least $2^n$ semantically distinct certified candidate compositions on that interface. If each such composition is attached by a certificate-preserving retrieval map to a declared problem descriptor, the corresponding declared problem family has at least $2^n$ certified candidate solutions.
\end{theorem}

\begin{proof}
There are $2^n$ subsets of $[n]$. By assumption each subset determines a certified admissible composition. Nondegeneracy makes the resulting semantic classes distinct, so the library has at least $2^n$ distinct candidate classes. A certificate-preserving retrieval map associates each declared descriptor with the corresponding certified class, yielding the second statement.
\end{proof}

\begin{definition}[Coherent realized task family and certified update work]
For a reusable family $M_n$, a coherent realized task family is a nested family of declared material task descriptors $T_n$ together with injective assignments
\begin{equation}
\Xi_n : \mathcal{P}([n]) \longrightarrow \{(\pi_I, p_I, \chi_I) : \pi_I \in T_n, p_I \text{ solves } \pi_I, \chi_I \in \mathrm{Cert}\}, \qquad \Xi_n|_{\mathcal{P}([n-1])} = \Xi_{n-1}, \tag{S38a}
\end{equation}
where $p_I$ is the certified composition associated with $I$, and $\chi_I$ is a target-side certificate checkable in the realization. Let $u_n > 0$ denote the amortized certified update work of adjoining $m_n$: module intake, index maintenance, composition-rule maintenance, and the certificate machinery that permits target-side verification on retrieval. Let $a_n$ be the number of newly activated descriptors from $\Xi_n(\mathcal{P}([n]) \setminus \mathcal{P}([n-1]))$ that are presented and resolved during the $n$-th update interval. The declared realized solving throughput is $R_n = a_n / u_n$.
\end{definition}

\begin{theorem}[Conditional certified solving-rate acceleration]
Let $L_D$ satisfy Theorem 12.25 and admit a coherent realized task family in the sense of Definition 12.26. Suppose that there exist $\alpha > 0$, $K > 0$, $N_0 \in \mathbb{N}$, and $1 \leq q < 2$ such that, for every $n \geq N_0$,
\begin{enumerate}[label=(R\arabic*)]
\item[(R1)] at least an $\alpha$-fraction of the newly indexed task-solution pairs is activated and target-side certified during the $n$-th update interval, so $a_n \geq \alpha 2^{n-1}$; and
\item[(R2)] the recorded amortized certified update work obeys the stepwise bound $u_{n+1} \leq q u_n$ and the envelope $u_n \leq K q^n$.
\end{enumerate}
Then the realized solving throughput has the accelerating certified lower bound
\begin{equation}
R_n \geq \underline{R}_n := \frac{\alpha 2^{n-1}}{u_n} \geq \frac{\alpha}{2K}\Big(\frac{2}{q}\Big)^n, \tag{S38b}
\end{equation}
and the lower-bound throughput satisfies
\begin{equation}
\frac{\underline{R}_{n+1}}{\underline{R}_n} = 2\frac{u_n}{u_{n+1}} \geq \frac{2}{q} > 1. \tag{S38c}
\end{equation}
Thus, in every realization satisfying (R1)--(R2), accelerated real task resolution is a globally coherent and empirically measurable consequence of the certified compositional architecture.
\end{theorem}

\begin{proof}
Coherence of $\Xi_n$ makes the $2^{n-1}$ subsets newly containing $m_n$ distinct from all prior indexed pairs. Injectivity and target-side certification make them distinct certified solutions of declared realized tasks. By (R1), at least $\alpha 2^{n-1}$ of those pairs are activated and resolved, so $a_n \geq \alpha 2^{n-1}$. Dividing by $u_n$ gives the first inequality in (S38b); (R2) gives the second. Finally,
\[
\frac{\underline{R}_{n+1}}{\underline{R}_n} = \frac{2^n / u_{n+1}}{2^{n-1}/u_n} = 2\frac{u_n}{u_{n+1}} \geq \frac{2}{q} > 1.
\]
The quantities $a_n$ and $u_n$ are declared operational measurements, so the consequence is empirically testable in any such realization.
\end{proof}

\begin{proposition}[Rate realization and empirical discharge]
Theorem 12.27 identifies the exact realization data for a measured solving-rate claim: a coherent real task family, activated certified descriptors, target-side verification, and recorded amortized update work. A task stream that realizes (R1) together with an implementation whose $u_{n+1}/u_n \leq q < 2$ realizes the theorem's acceleration prediction; measured departures localize to task activation, semantic composability, certificate preservation, or update-work growth.
\end{proposition}

\begin{proof}
Each named datum occurs in Definition 12.26 or premises (R1)--(R2). Theorem 12.27 gives the resulting lower bound. Failure of an implementation to realize the conclusion therefore identifies a failed or unverified premise rather than an untyped gap between capacity and rate.
\end{proof}

\begin{definition}[Reference implementation contract]
A realized semantic library records (i) semantic canonicalization or a typed approximation residual, (ii) certificate checks, (iii) task-descriptor activation logs, (iv) separated meters for composition, retrieval, verification, and update work, and (v) a preregistered report with three outcomes: realized under the declared premises, named-premise failure, or unresolved within the declared budget. This is an implementation obligation, not an additional premise of the theorems above.
\end{definition}

\begin{remark}[Dyadic library witness]
The realization class is nonempty. Take modules indexed by binary places, let each subset compose to the corresponding binary sum, use exact canonical binary representation as the semantic key, attach an exact certificate to each sum, and activate the descriptors indexed by those sums. Distinct subsets have distinct binary sums, so the $2^n$ candidate family is realized exactly. Choosing a declared update meter with $u_{n+1} \leq q u_n$ for some $q < 2$ realizes the algebraic premises of Theorem 12.27. This finite witness fixes an audit template for the realization conditions without altering the theorem's positive premise-complete statement.
\end{remark}

\begin{remark}[Topological-map reading]
The library is a topological map of a functional state space: vertices are semantic classes of conditions or certified subsolutions, and edges are admissible transformations with declared evaluation consequences. A well-formed problem is the absence of a known certified path between typed endpoints; a solution is a path composition. The admissibility, equivalence, and certificate layers are what make this map more than an undifferentiated association network --- and the reading requires no geometric drawing and no completed global conceptual manifold: it is exact on the library as built.
\end{remark}

\noindent\plainterms{} Three results now form one cumulative mechanism. First: filed by meaning and updated only through certificates, the shelf of solved routes never gets worse. Second: hold $n$ compatible certified building blocks and you hold $2^n$ distinct certified assemblies, provided the six conditions hold. Third: when those assemblies are coherently attached to distinct real tasks, a nonzero share of those tasks is activated, and the full certificate-accounted work of adding a module grows by less than a factor of two per step, newly certified real task resolutions per unit of update work accelerate. Retrieval and verification are not waved away: they are inside $u_n$, the quantity the theorem prices and measures. The theorem therefore makes a positive conditional prediction about real problem-solving rate, with the exact conditions under which an implementation realizes it.

\subsection{Relation to the Cognitive Near-Singularity}

\begin{remark}[CNS relation: substrate and convergence regime, registered]
A persistent, compositional, certificate-preserving semantic library is a material substrate condition for a Cognitive Near-Singularity regime: solved paths become reusable operands in later path construction, supplying the growth material that a near-singular law consumes. The published CNS possibility result supplies its own contraction, composition-closure, invariance, and witness conditions [13]. These conditions are also a candidate realization regime for the convergent side of recursive HCFM assessment: where the independent CNS witness conditions and the assessment conditions of Theorem 12.17 are jointly instantiated, correction can be governed by a common fixed-point structure rather than only accumulate. The present rate theorem supplies the complementary cumulative-production route. This is a T5 structural relation: each contribution retains its own premises, witness, and realization conditions.
\end{remark}

\subsection{Reference governance and module ledger}

This subsection is the single home of the module's status bookkeeping: the tag vocabulary, the dependency records through which narrower papers are projected from this reference, and the promotion conditions under which each architectural claim becomes a stronger one.

\begin{definition}[Reference-status tag]
Every item in this module and the alignment module carries one of the following status tags.
\begin{enumerate}[label=(T\arabic*)]
\item[(T1)] Construction: follows from a prior proved construction by the stated definition or proof.
\item[(T2)] Definition: fixes a typed object or protocol and makes no existence claim beyond its declared input data.
\item[(T3)] Conditional theorem: derives a conclusion only from displayed premises.
\item[(T4)] Implementation obligation: identifies a datum, bridge, algorithm, or certificate that a realizing system must provide.
\item[(T5)] Frontier relation: records a connection to another theory without using it as a premise or consequence here.
\end{enumerate}
A canonical reference is adequate when a reader can recover each item's status, inputs, output type, and promotion or failure condition without reconstructing them from prose.
\end{definition}

\begin{remark}[The conceptual-space interface]
The phrase ``human conceptual space'' is used in this reference conditionally: if a globally coherent conceptual space $K$ is supplied, a discipline is represented through a typed interface into $K$, and this module's results state exactly what that interface must preserve. The results are proved for discipline-local interfaces, so they hold now and compose with a global $K$ whenever one is constructed, partially realized, or supplied by another program.
\end{remark}

\begin{definition}[Canonical-reference dependency record]
For each nontrivial disciplinary claim $q$, a dependency record is the tuple
\begin{equation}
\mathrm{DepRec}(q) = \big(\mathrm{type}(q), \mathrm{inputs}(q), \mathrm{status}(q), \mathrm{scope}(q), \mathrm{certificate}(q), \mathrm{failure}(q), \mathrm{projection}(q)\big), \tag{S39}
\end{equation}
recording its input types, status tag, valid scope, needed certificate, condition that blocks promotion, and the narrower carrier forms that may later cite or specialize it. A reference document is derivation-ready for a claim family when every load-bearing claim has such a record and no proof cites a claim whose status is weaker than the proof requires.
\end{definition}

\begin{center}
\small
\begin{longtable}{@{}p{0.24\textwidth}p{0.08\textwidth}p{0.30\textwidth}p{0.28\textwidth}@{}}
\toprule
Component & Status & Dependencies / certificate & Additional condition for the stronger claim\\
\midrule
\endhead
Disciplinary realization (S29) & T2 & Typed domain, operators, evaluator, equivalence, certificates, perspective & A constructed universal conceptual space, for the global-embedding claim\\
No-silent-transport proposition & T1 & Source/target type check & A constructed bridge, for any particular transport claim\\
Perspective-explicit bridge & T2 & Clauses (B1)--(B6) & Certified clause discharge, for evaluator and certificate preservation\\
Disciplinary invariant transport & T1 & Certified bridge plus dynamic equivariance & A certified bridge on the wider domain, for scope extension\\
Three-valued coherence assessment & T2 & Declared HCFM system model, perspective, budget, noise floor, ceiling & A coherence witness or valid counter-witness for a non-undetermined verdict\\
Observer-inclusion blind-spot proposition & T3 & Unbounded certified corrections, omitted assessment condition, and persistent residual above noise & A realized recursively extensible assessment class with measured residuals\\
Recursive global-coherence convergence theorem & T3 & A1--A6: fairness, soundness, finite separation, completion, convergence, decisive resolution & A realized assessment scheduler, verifier, and completion certificate\\
Semantic backpropagation & T2 & Quotient-factored evaluator, library graph, budget & Retrieval indexing and benchmarking, for implementation-specific gain estimates\\
Retention monotonicity & T1 & Certified update rule and finite bounded candidate class & Coherent task realization and update-work accounting, for a throughput conclusion\\
Candidate capacity and rate theorems & T3 & C1--C6, nondegenerate modules, coherent task realization, and (R1)--(R2) & Empirical instantiation of task activation and recorded update work\\
CNS relation & T5 & Separate CNS conditions, witness, and Theorem 12.27 realization conditions & Joint realization of the independent CNS premises\\
\bottomrule
\end{longtable}
\end{center}

\begin{obligation}[Audit protocol for a derived carrier]
Before a narrower paper, implementation, or empirical proposal is projected from this reference, audit: (i) the exact source claims used; (ii) their status tags; (iii) all omitted input conditions; (iv) the target discipline and bridge scope; (v) the certificate that survives the projection; and (vi) the new falsifier introduced by the target carrier. A projection may simplify a claim only when it preserves this record; it may not promote a definition, conditional theorem, or frontier relation by omission.
\end{obligation}

\begin{center}
\small
\begin{longtable}{@{}p{0.26\textwidth}p{0.30\textwidth}p{0.36\textwidth}@{}}
\toprule
Candidate derived carrier & Reference projection & Additional work before independent claim\\
\midrule
\endhead
Second-order disciplinary methodology & Definitions 12.1--12.5 and Proposition 12.3 & Exemplary discipline pairs and bridge-audit cases\\
Coherence-audit protocol & Definitions 12.8--12.16 and Theorem 12.17 & Domain-specific residual scales, fair scheduler, verifier, and completion or refutation certificates\\
Semantic-backpropagation system & Definitions 12.19--12.20 and Proposition 12.21 & Canonicalization, evaluator implementation, benchmarked certificate checks\\
Compositional problem-solving architecture & Definition 12.24 and Theorems 12.25, 12.27 & Realized task family, activated descriptors, retrieval index, verifier, and recorded update-work evaluation\\
CNS-oriented systems analysis & Remark 12.32 & Independent proof or witness of the CNS conditions for the constructed trajectory\\
Cross-disciplinary invariant-transport study & Definition 12.5 and Proposition 12.6 & Domain-specific dynamic-equivariance and certificate-translation evidence\\
\bottomrule
\end{longtable}
\end{center}

\begin{remark}[Promotion conditions]
Each architectural claim of this module states, in its dependency-table row, the exact condition under which its positive consequence is realized: a completed conceptual-space embedding requires a constructed $K$; decidable semantic equivalence requires a decision procedure for the declared $\equiv_D$; combinatorial capacity requires (C5)--(C6) on top of (C1)--(C4); certified solving-rate acceleration requires the coherent task realization and update-work conditions of Theorem 12.27; and a CNS conclusion requires independent discharge of the CNS conditions. These are the module's realization ladder: every rung is named, and nothing on a lower rung borrows the authority of a higher one.
\end{remark}

\section{Alignment dynamics: windows, cognitive mass, and realization gating}

The construction says what a realized observer model is; this module records what the same machinery implies about how such models grow, reorganize, or fail --- and where control is possible. It is a separately typed hypothesis layer, and the ledger in \S13.8 is the authoritative record of each item's status.

\noindent\textbf{Type discipline.} Items below use the status vocabulary of Definition 12.33: definitions are (T2); propositions carrying proofs from displayed premises are (T3); registered conjectures and the policy schema carry explicit falsifiers and dependencies and are used in no proof; obligations are (T4). Throughout, $(C, H) \in \ObsMatcl$ satisfies the observer premises of Theorem 22, and order counts meta-lift depth in the sense of Definition 38 and Theorem 39.

\subsection{Operational windows and the boundary dichotomy}

\begin{definition}[Operational window]
For an HCFM realization at order $n$, the operational window is
\[
W_n(H) = [\varepsilon_n(H), K_n(H)],
\]
where $\varepsilon_n(H)$ is the declared resolution of $H$ (the noise floor: the finest operational distinction separated by the protocol family), and $K_n(H)$ is the observability ceiling: the supremum of composition scales $s$ --- path length in the registration path category, weighted by the declared registration-label cost of Definition 11.1 --- such that the readout of compositions of scale $s$ remains within the observer's protocol family, at the declared resolution, over the viable accessed region $V_{C,H}$. The accessed scale $s(t) \in W_n(H)$ is the scale of the composition class exercised by the system's admissible continuation at time $t$. Compression is monotone motion of $s(t)$ toward $K_n$; expansion is monotone motion toward $\varepsilon_n$. (T2)
\end{definition}

\begin{remark}[Window and scale ladder]
The operational window $W_n(H)$ is a floor-and-ceiling dynamic range; the scale ladder $W_H$ of Definition 11.1 is the set of measuring scales at which $d_A$ is evaluated. The two instruments compose --- $d_A$ is evaluated inside $W_n$, and a horizon whose ladder cannot certify a covering exponent reports $\bot$ --- and each keeps its own name. (T2)
\end{remark}

\noindent\plainterms{} Every working system --- a mind, an instrument, an institution, a model --- has a finest grain it can resolve and a largest load it can hold together at once. The stretch between those limits is its operational window: the band in which it actually functions. ``Compression'' means its workload is climbing toward the ceiling; ``expansion'' means its distinctions are thinning toward the floor. (The scale ladder is a different tool --- a set of measuring rings for dimension --- and keeps its own name so the two are never confused.)

\begin{definition}[Coherent order transition; decoherence]
An encounter of $s(t)$ with $\partial W_n(H)$ has a coherent continuation when there exists a faithful, flow-derived encoding-and-lift scheme (Definition 38) carrying the realization to order $n+1$ with window $W_{n+1}(H')$, such that the boundary compositions of order $n$ are primitives (atoms of the typed presentation) at order $n+1$ and the lift preserves the domain types, dynamic bases, path composition, bridges, and fitness objects (Theorem 39). The encounter ends in decoherence when no such lift exists and the operational quotient fails to remain well-defined under continued dynamics: the greatest horizon bisimulation of Theorem 16 ceases to be stable under the flow on the accessed region, or the strict adaptive contraction of Theorem 22 is lost, so the structure exits the class of functional state spaces at the declared horizon. (T2)
\end{definition}

\begin{conjecture}[Boundary dichotomy]
For systems satisfying the premises of Theorem 22, every driven boundary encounter (one in which $s(t)$ reaches $\partial W_n$ under nonzero influx in the sense of Definition 13.16) terminates in exactly one of the two continuations of Definition 13.3; no third stable continuation exists. Registered conjecture; used in no proof; cited as an explicit premise by Proposition 13.18.
\end{conjecture}

\noindent\plainterms{} At the edge of its window, a driven system either reorganizes --- yesterday's whole procedures become today's single building blocks, and the window itself is re-based one level up --- or it comes apart: it can no longer keep its own states distinct or correct its own drift. Reorganize or unravel. That there is no third stable outcome is stated as a conjecture, and everything downstream that leans on it says so out loud.

\begin{remark}[Readings and scope]
The physical reading of the compression branch is the gravitational-collapse ladder: saturation of a degeneracy floor, re-basing to deeper primitives where a coherent phase exists, collapse where none does; each threshold crossing is driven by matter external to the compressing core. The cognitive reading is the transition between reasoning orders, in which whole order-$n$ compositions become order-$(n+1)$ primitives. The expansion branch is anchored cognitively (generalization either re-coheres as a coarser faithful quotient or dissolves as distinctions approach the floor) and is scoped to bound subsystems with external drivers; cosmological expansion lacks the external-driver structure assumed throughout this module and lies outside its scope. (T2/T5)
\end{remark}

\subsection{Registered coupling and the signed mass invariant}

\begin{definition}[Registered coupling; signed cognitive mass]
Let $X$ be a structure external to $S = (C, H)$, and let $c \geq 0$ parametrize registered coupling of $S$ to $X$: the component of interaction separated by $S$'s protocol family above $\varepsilon_n$. Write $\mathrm{Reach}_H(S; c)$ for a declared reach functional on the reachable operational quotient within $W_n(H)$. The signed mass of $X$ relative to $S$ at $H$ is
\begin{equation}
m_S(X; H) = \frac{\Delta \mathrm{Reach}_H(S; c)}{\Delta c}, \tag{S45}
\end{equation}
evaluated on a declared finite-difference domain. Coupling below the separating scale has registered component and mass zero. Positive mass expands registered reach; negative mass constrains or coordinates it. (T2)
\end{definition}

\begin{proposition}[Signed mass is an observable invariant]
$m_S(X; H)$ is constant on operational equivalence classes and factors through the registration quotient. It is therefore an order-relative observable invariant. A coherent order lift can make previously unresolved coupling mass-visible; both signs are admissible. (T3)
\end{proposition}

\begin{proof}
Registered coupling and reach are quotient-level objects, so the finite difference is class-constant and factors through the quotient by Proposition 28. Below the separating scale the registered increment is zero. A faithful lift can separate a formerly unresolved coupling. The sign is the direction of the reach change; no nonnegativity premise is required.
\end{proof}

\begin{proposition}[Order-relativity of signed mass]
If the least separating scale of $X$'s interaction with $S$ exceeds $K_n(H)$, then $m_S(X; H) = 0$ at order $n$. If a coherent transition to order $n+1$ raises the ceiling above that scale, the lifted quotient can assign nonzero signed mass whenever the newly registered coupling changes reach. The same external structure can therefore have different signed mass at different orders. (T3)
\end{proposition}

\begin{proof}
The registered coupling increment vanishes below the window. After a lift satisfying Theorem 39, the lifted protocol family can separate the interaction and the declared reach functional may change. The sign follows from Definition 13.6.
\end{proof}

\noindent\plainterms{} ``Mass'' is signed reach change per unit of input a system can actually register. A source can open reachable possibilities or constrain them into a coordinated form; both effects are observable only relative to the observer's current separating scale. A capability transition can make a formerly invisible source mass-visible without changing the source itself.

\begin{remark}[Starvation is coupling reduction]
Every implementable ``starvation'' intervention acts on registered coupling $c$, not on the substance $X$: compartmentalization, decoupling, and unregistrability change the coupling relation while leaving the source present. The gravitational correspondence is heuristic and typed as such (T5). This motivates the load definition below.
\end{remark}

\subsection{Nucleation}

\begin{conjecture}[Nucleation requirement]
There exists $m_{\mathrm{crit}}(n) > 0$ such that a boundary encounter at $K_n(H)$ has a coherent continuation only if the total registered external mass satisfies $\sum_X m_S(X; H) \geq m_{\mathrm{crit}}(n)$, the sum ranging over external structures independent of $S$ --- including independent instances of $S$ itself, but excluding $S$'s own outputs re-registered as inputs. A system driven to its ceiling with total independent external mass below threshold decoheres. Registered conjecture; falsifier in Obligation 13.12.
\end{conjecture}

\begin{remark}[Anchors and the self-play case]
Developmental anchor: order transitions in human cognition empirically require external scaffolding slightly ahead of current capacity. Physical anchor: collapse-ladder thresholds are crossed under external load. The apparent machine-learning counterexample of self-improvement by self-play conforms on inspection: self-play manufactures an independent peer instance plus fixed external rule structure --- nucleation by cloning. The exclusion clause in Conjecture 13.10 is therefore essential to its content: isolation means absence of any independent registrable structure, including copies. (T5)
\end{remark}

\begin{obligation}[Falsifier for nucleation]
One verified coherent transition to order $n+1$ under confirmed coupling isolation (in the sense of the exclusion clause) refutes Conjecture 13.10. Upon refutation, Policy Schema 13.20 loses its rationale and Conjecture 13.22 is withdrawn; this dependency must remain explicit wherever those items are cited. (T4)
\end{obligation}

\noindent\plainterms{} The claim: the leap to a new level of organization cannot be powered purely from inside. Something independent --- a teacher slightly ahead, a peer, a rival, an environment with its own structure --- must seed the new pattern, the way a supersaturated solution crystallizes around a speck and, without one, just curdles. The claim is falsifiable and says how: exhibit one system that verifiably made the leap in genuine isolation (no independent structure, not even a copy of itself to push against), and this conjecture --- and the safety policy built on it --- falls.

\subsection{The two-channel decomposition}

\begin{definition}[Assessment and realization coupling]
Let $\chi$ be a composition held in the conceptual/mathematical columns of a source realization. Coupling of an external system $S'$ to $\chi$ is assessment coupling when it is mediated by registration of the formal object itself, through the conceptual and mathematical bridges of Definition 19; this requires $\chi$'s presentation scale to lie within $W_{n'}(H')$ of the registering system --- an order-gated condition. A material-realization interface for $\chi$ is declared realization enrichment that carries the Real-interpreted definable object of $\chi$ to a physical implementation and consequence trace, with its own certificates and residuals. Coupling is realization coupling when it is mediated by material consequences through such a declared interface; this requires of $S'$ only material coupling and base-order registrability of the consequences. (T2)
\end{definition}

\begin{proposition}[Population gating]
Fix a population $P$ of coupled observer systems and a composition $\chi$ of presentation scale $\sigma$. The total mass of $\chi$ over $P$ decomposes as
\[
M_P(\chi) = \underbrace{\sum_{S' \in P: \sigma \in W(S')} m^a_{S'}(\chi)}_{\text{assessment channel}} + [\chi \text{ realized}] \cdot \underbrace{\sum_{S' \in P} m^r_{S'}(\chi)}_{\text{realization channel}},
\]
where the first sum ranges only over the order-capable subpopulation and the second over all materially coupled members. Realization changes the registrable population discontinuously from the first index set to the second. (T3)
\end{proposition}

\begin{proof}
The two summands instantiate Definition 13.6 on the two coupling routes of Definition 13.13; the index sets are forced by the registrability conditions of the respective columns (order-gated presentation for C/M registration; base-order material consequence for the physical column). The bracket is the realization indicator: before realization the second channel's registered coupling is zero by definition.
\end{proof}

\begin{remark}[Empirical discharge as the mass event]
Proposition 13.14 is the formal content of the informal observation that effective cognitive mass tracks local empirical demonstration more strongly than global coherence assessment: not because assessment is epistemically lesser, but because its registrable population at any moment is the small order-capable subpopulation, whereas realization couples to every materially connected observer. The demonstration, not the theorem, is the mass injection. (T3 gloss)
\end{remark}

\noindent\plainterms{} A blueprint moves only those who can read blueprints; the running machine moves everyone who sees it run. The proposition just does the accounting: total impact = (impact per capable reader) $\times$ (number of capable readers) +, once built, (impact per witness) $\times$ (everyone). Building is the step that flips a result from the narrow channel into the universal one --- which is why the moment of maximum inflow into the world is a working demonstration, not a proof.

\subsection{Influx, horizon crossing, and the intervention window}

\begin{definition}[Load influx, horizon crossing, and intervention window]
For a system or collective $S$, write $i(t) = e(t) + g(t)$ for registered load-influx, weighted by $|m|$: $e(t)$ is externally gateable load and $g(t)$ is load generated when realized outputs are re-registered across members. Write $\kappa(t)$ for coherence-capacity growth. The nucleation condition declares whether its threshold consumes expansive mass $m_+$ or total organizing/expansive load $|m|$; this choice is an explicit falsifiable parameter. $S$ is horizon-crossed at $t^{*}$ when internal $g(t)$ remains sufficient to sustain the declared drive law without external $e$. The intervention window is $[t_{\mathrm{now}}, t^{*})$: gating acts on $e$, while capacity growth changes $\kappa$.
\end{definition}

\begin{definition}[Drive-coupled realization]
Write $x(t) = s(t)/K_n(t) \in (0, 1]$ for relative accessed scale and $\kappa(t) = \tfrac{d}{dt} \log K_n(t)$. A realization is drive-coupled with parameters $\beta > 0$, $\gamma \geq 0$ when, outside deliberately initiated transitions,
\begin{equation}
\frac{d}{dt} \log x(t) \geq \beta i(t) - \gamma - \kappa(t). \tag{S47}
\end{equation}
A trajectory for which the integral of the displayed right-hand side diverges while $x$ stays bounded away from one falsifies drive coupling for that realization.
\end{definition}

\begin{lemma}[Finite-time boundary encounter]
Under (S47), if $\beta i(t) - \gamma - \kappa(t) \geq \delta > 0$ on $[t_0, t_1]$, then $x(t)$ reaches one by time $t_0 - \log x(t_0)/\delta$, unless the trajectory exits the window earlier.
\end{lemma}

\begin{proof}
Integrating (S47) yields $\log x(t) \geq \log x(t_0) + \delta(t - t_0)$; the stated hitting bound follows.
\end{proof}

\begin{proposition}[Two-lever control claim]
Assume the boundary dichotomy, the declared nucleation requirement, and drive coupling (S47). A never-decohering trajectory either keeps the safety budget $\beta i(t) \leq \kappa(t) + \gamma$ outside deliberately nucleated encounters or enters each supercritical encounter nucleation-ready. Influx throttling reduces the external $e$-component of $i$; capacity growth raises $\kappa$ and therefore enlarges the admissible load budget. The claim is conditional on these displayed premises and is not used by the material-generative core.
\end{proposition}

\begin{proof}
Persistent violation of the budget gives the positive margin of Lemma 41, hence a driven boundary encounter. The boundary dichotomy and nucleation condition determine the permitted coherent continuation. The two lever roles follow from their displayed terms in (S47).
\end{proof}

\begin{remark}[Named failure modes]
Throttling-only fails by defection instability ($e$-control is multilateral; any defector restores influx) and by the bootstrapping paradox (reliable collective-scale gating presupposes coordination capacity of the order the gate protects against). Capacity-only fails under shock ($\kappa$ is slow relative to influx spikes). Viable trajectories alternate: throttle to buy time, raise the ceiling, then admit mass and transition deliberately. (T3 gloss)
\end{remark}

\noindent\plainterms{} Two dials. One slows what comes in: instant, but it needs everyone's cooperation forever, and one defector spins it back --- worse, running such a gate over a whole civilization already requires roughly the coordination capacity the gate is supposed to protect. The other grows what the system can digest: slow, but it compounds and one actor can turn it alone. All-throttle freezes you below a leap you will eventually need; all-capacity is too slow for shocks. So the conditional claim: viable paths alternate --- slow the inflow, grow the digestion, then let it in on purpose. And the window for the first dial is finite: once a system feeds mainly on its own outputs, outside throttles stop reaching it at all.

\subsection{The realization gate}

\begin{policyschema}[Realization gate]
Let $\Theta$ be a declared collective coherence threshold for a population $P$: the complexity scale above which no individual or institutional member of $P$ achieves the declared coherence conditions unaided ($\Theta$ is a declared datum of the schema instance, of the same status as the tolerance $\tau_H$ in Definition 11.1). A composition $\chi$ is admissible for material realization through a declared material-realization interface in the sense of Definition 13.13 if and only if:
\begin{enumerate}[label=(G\arabic*)]
\item[(G1)] there exists a problem $\Pi$ of complexity above $\Theta$ such that every admissible continuation of $P$ in which $\Pi$ is unsolved exits all operational windows or violates viability (existential necessity, typed);
\item[(G2)] $\chi$ belongs to a minimal set of compositions jointly sufficient to solve $\Pi$; and
\item[(G3)] assessment of $\chi$ --- including simulation, which is assessment extended by computational means --- has been carried to the assessment-fidelity horizon of Remark 13.21 without producing a coherence failure, with the outcome recorded as a certificate in the sense of Definition 12.19 and entered in the dependency record of Definition 12.35.
\end{enumerate}
Assessment coupling to $\chi$ is unrestricted under this schema. Type: normative policy schema, conditional on Conjectures 13.4 and 13.10 and on Proposition 13.14. Enforcement is protocol-internal: the gate is an admissibility condition inside the certificate discipline --- what certified processes will compose with --- since a decentralized collective admits no external throttle.
\end{policyschema}

\begin{remark}[Assessment-fidelity horizon]
As the scale of $\chi$ approaches the collective ceiling, the cost and coupling of faithful simulation of $\chi$ approach the cost and coupling of realizing it, and unfaithful simulation gates nothing. Gating-by-simulation therefore carries an intrinsic horizon, and the binding constraint on Policy Schema 13.20 is simulation capacity. Machine-checkable kernels, typed-constraint verifiers, and conformance suites are correctly typed as horizon extension for clause (G3). The same accounting is explicit in Theorem 12.27: certificate maintenance and verification belong inside the recorded work term $u_n$, so improved verifier infrastructure directly improves the realizable throughput regime. (T3 gloss)
\end{remark}

\noindent\plainterms{} The gate, in one sentence: think, model, and simulate anything; build only what an existential problem demands, only the minimal set that solves it, and only after simulation-to-the-limit has been recorded with a certificate. Two design facts follow from the mechanics rather than from preference. The rule must live inside the protocols --- in what certified work agrees to compose with --- because nobody stands outside a decentralized crowd holding a valve. And the rule is only as strong as your verification tools, because near the limit a simulation faithful enough to trust costs nearly as much as the real thing; that is why verifier infrastructure is part of the gate, not an accessory to it.

\subsection{Implications, selection reading, and the preregistered sign question}

\begin{conjecture}[Inward decoherence; anti-foom]
Assume Conjecture 13.10. A self-coupled system whose only influx is $g(t)$ --- no independent external structure in the sense of the exclusion clause --- and which is driven to its ceiling does not undergo a coherent order transition; it decoheres. Sustained capability compression through successive orders is therefore available only to systems maintaining registered coupling to independent external structure. Dependency: withdrawn if Obligation 13.12 fires.
\end{conjecture}

\begin{definition}[Preregistered sign question]
For diffusion of order-$(n+1)$ capacity through a population $P$, define $R_{\mathrm{cap}}(t) = \dot{\Theta}(t)$ (collective capacity growth attributable to diffusion) and $R_{\mathrm{mass}}(t) = \dot{g}(t) + \dot{e}(t)$ (collective mass-generation growth attributable to diffusion; member realizations enter other members' $e$, scaling superlinearly in $|P|$). Preregistered outcomes: (O1) $R_{\mathrm{cap}} > R_{\mathrm{mass}}$ robustly across regimes --- diffusion is self-protective; the capacity lever dominates and gating is transitional. (O2) $R_{\mathrm{mass}} > R_{\mathrm{cap}}$ --- diffusion is self-igniting; the gate binds and must precede diffusion-scale demonstrations. (O3) no stable sign --- control must be adaptive and regime-indexed. All three outcomes are informative and must be reported with equal prominence. Measurement path: define both rates via Definition 13.6 in bounded certified collectives of increasing size. (T2/T4)
\end{definition}

\begin{remark}[Great Filter typing]
The Great Filter reading of this module is population-level survivorship: if driven boundary encounters generically end in decoherence absent nucleation and gating (Conjectures 13.4 and 13.10), then long-lived observed systems are those that remained within windows or transitioned coherently. This is an observation-selection consequence of the mechanism --- the same inference form as the conditional dimensional-resolution result (Theorem 11.15), applied to persistence rather than to registered dimension --- and introduces no additional mechanism. The two readings of ``self-selecting'' (within-system fitness-following; population survivorship) each keep their own role: the first supplies dynamics, the second supplies the Filter. (T5)
\end{remark}

\begin{remark}[Operational corollary: proof before demonstration]
By Proposition 13.14, a working demonstration of a high-impact collective process is a realization-channel mass injection over the full coupled population, whereas formal dissemination remains in the order-gated assessment channel. Under Policy Schema 13.20, certification and gating infrastructure must therefore precede demonstration for any composition whose realization scale approaches the collective ceiling; diffusion sequencing that leads with demonstration is, on this module's own terms, the horizon-crossing event. (T3 gloss, conditional as its premises)
\end{remark}

\begin{remark}[Witness tie-in]
The coupled-clock witness (Section 14) is a microscale illustration of Definition 13.6: the emitted pair interaction is registered coupling between the clocks, and the computation there shows it reshaping the composite's viable operational quotient --- reducing effective registration dimension from the uncoupled product value to one --- so that what each subsystem can reach and distinguish is altered by coupling alone, with nothing added or removed. The witness also exercises the $\bot$ clause, the microscale analogue of a window too coarse to certify its own structure. The tie is illustrative (T5); the module's claims stand on their own stated premises.
\end{remark}

\noindent\plainterms{} Three consequences with their labels showing, then the open question. (1) The isolated self-improver is predicted to unravel, not explode --- leaps need outside seeds; falsifier stated above. (2) The mirror-image danger: in a crowd of builders, everything anyone builds feeds everyone else, so inflow grows faster than headcount. (3) Old observed systems are filter-survivors --- a statement about what we get to see, in exactly the pattern of the earlier three-dimensions result. The open question is the sign: does spreading capability grow the crowd's digestion faster than its appetite? The module commits in advance to how that would be measured and to reporting whichever answer arrives. And one sequencing rule falls out before any measurement: for anything big, the certificates and the gate come before the public demo --- because the demo is the flood.

\subsection{Module ledger and scope}

\begin{center}
\small
\begin{longtable}{@{}p{0.36\textwidth}p{0.20\textwidth}p{0.38\textwidth}@{}}
\toprule
Item & Type & State at inclusion\\
\midrule
\endhead
Def. 13.1, Rem. 13.2 (window, ladder) & Definition (T2) & Complete; type-separated from $W_H$\\
Def. 13.3 (transition/decoherence) & Definition (T2) & Complete; consumes Def. 38, Thms. 39, 16, 22\\
Conj. 13.4 (dichotomy) & Conjecture & Used in no proof; explicit premise of Prop. 13.18\\
Def. 13.6 (mass) & Definition (T2) & Complete\\
Prop. 13.7 (invariance) & Proposition (T3) & Proved via Prop. 28\\
Prop. 13.8 (order-relativity) & Proposition (T3) & Proved via Thm. 39\\
Conj. 13.10 (nucleation) & Conjecture & Falsifier named (Obl. 13.12)\\
Def. 13.13, Prop. 13.14 (two channels) & Definition; Proposition (T3) & Proved from registrability conditions\\
Def. 13.16 (influx, crossing, window) & Definition (T2) & Complete\\
Prop. 13.18 (two levers) & Conditional (T3) & Conditional exactly on the two conjectures\\
Policy Schema 13.20 (realization gate) & Normative schema & Conditional; enforcement protocol-internal via Cert and Def. 12.35\\
Rem. 13.21 (fidelity horizon) & Remark (T3 gloss) & Complete; verifier infrastructure typed as horizon extension\\
Conj. 13.22 (anti-foom) & Conjecture & Dependency on nucleation explicit\\
Def. 13.23 (sign question) & Preregistration & O1--O3 typed; measurement path named\\
Rem. 13.24, 13.25, 13.26 & Remarks & Complete\\
\bottomrule
\end{longtable}
\end{center}

Realization scope of this module, recorded once: quantitative values of $m_{\mathrm{crit}}(n)$, $\theta_{\mathrm{drive}}$, and $\Theta$, the sign-question measurement, historical retrodiction case studies, and any cosmological application are empirical realization data for particular domains. Under the module's declared transition and nucleation premises, the control, survivorship, and sequencing predictions are globally coherent; their measurement identifies which realized regime obtains. Dependency declaration: no item in this module is used in the proof of any result in the preceding sections; Theorem 40 and its proof are unchanged by this module, which consumes the dimensionality and disciplinary modules only as data.

\section{Coupled-clock witness}

The witness executes the whole construction on a material system whose interaction, quotient, invariant, dimensional report, carrier, and bridge can all be computed.

\begin{remark}[Witness scope]
The coupled-clock witness establishes joint satisfiability and end-to-end computability of the declared constructions on its stated viable quotient. It provides a possibility witness for interaction-dependent registration structure, including the honest unresolved outcome when the ladder fails. Its scope is the declared oscillator system and horizon; generic rank-three realization, universal interaction-induced dimension reduction, and the alignment conjectures retain their independent conditions.
\end{remark}

\begin{remark}[Detuned lock/drift extension and independent anchor]
For two detuned phase oscillators $\dot{\theta}_1 = \omega_1 + K \sin(\theta_2 - \theta_1)$, $\dot{\theta}_2 = \omega_2 + K \sin(\theta_1 - \theta_2)$, the classical Adler/Kuramoto lock--drift boundary separates a phase-locked regime from a drifting regime [33, 34, 35, 36]. Under the declared finite-resolution registration ladder, locking yields a one-dimensional collective-phase quotient and drifting yields a two-dimensional accessible phase quotient; near threshold, finite resources can leave the covering profile unresolved. The classical boundary is prior art; the HCFM contribution is the registration-dimension reading and its finite-resolution prediction.
\end{remark}

\begin{remark}[Certificate-bearing bridge audit]
A worked bridge between two realized descriptions must record: (B1) translation of objects and operations; (B2) admissibility residue; (B3) value tolerance; (B4) dynamical commutation residual; (B5) certificate transport; and (B6) transported-invariant or round-trip residual. The coupled-clock phase-coordinate bridge satisfies these clauses when its coordinate map preserves the declared flow and readout bins; any residual is reported at the clause where it occurs. This audit is the operational form of perspective-explicit transport, not an optional narrative comparison.
\end{remark}

Let $C = C_1 \otimes C_2$ be two circle-valued clock components with angles $(\theta_1, \theta_2) \in T^2$, equal uncoupled frequency $\omega$, and material generator
\begin{equation}
X_C = \omega(\partial_{\theta_1} + \partial_{\theta_2}) + \kappa \sin(\theta_2 - \theta_1)(\partial_{\theta_1} - \partial_{\theta_2}). \tag{S24}
\end{equation}
The global law on the clock family is the action generated by compatible restrictions of this vector field to one-clock and two-clock composites. The singleton emitted components are $\omega \partial_{\theta_i}$, and M\"obius extraction yields the pair component
\begin{equation}
X_C(\{1, 2\}) = \kappa \sin(\theta_2 - \theta_1)(\partial_{\theta_1} - \partial_{\theta_2}). \tag{S25}
\end{equation}
This component is nonzero precisely when the composite differs from independent clock motion. Put $\delta = \theta_2 - \theta_1$ and $\phi = (\theta_1 + \theta_2)/2$. For $\kappa > 0$, (S24) gives $\dot{\delta} = -2\kappa \sin\delta$, so the phase-locked regime $\delta = 0$ is forward stable at the declared phase resolution.

The observer uses the same two phase probes in both comparisons and reads each phase in $N$ quantization bins. The registration label cost assigns unit cost to an adjacent resolvable phase move. Let the declared scale ladder contain radii $1 \leq r < R \leq N/8$ and a tolerance $\tau_H < 1/4$.

\noindent\textbf{Uncoupled comparison regime.} For $\kappa = 0$, take the viable accessed registration region to be the full quantized phase torus $C_N \Box C_N$, with its graph metric. The two phase probes are independently separating. Covering a graph ball at the declared scales has two-dimensional growth: there are constants independent of $r, R$ for which
\begin{equation}
a(R/r)^2 \leq N_V(x; R, r) \leq b(R/r)^2. \tag{S26}
\end{equation}
After choosing the ladder so that the corresponding logarithmic constants are below $\tau_H$, Definitions 11.1--11.2 give $d_A(H_{\mathrm{unc}}) = 2$. The same calculation for either single clock, whose quotient is the cycle graph $C_N$, gives $d_A(H_{C_1}) = d_A(H_{C_2}) = 1$, in agreement with Proposition 11.9.

\noindent\textbf{Coupled viable regime.} For $\kappa > 0$, take the viable observer region to consist of histories whose difference coordinate has entered and remains inside one declared resolution bin about $\delta = 0$. The probe family is unchanged, but the coupling-induced stable relation makes all residual $\delta$-variation operationally nonseparating at that resolution. The resulting registration quotient is the collective-phase cycle $C_N$ in the $\phi$-coordinate. Its covering profile is linear:
\begin{equation}
a'(R/r) \leq N_V(x; R, r) \leq b'(R/r), \tag{S27}
\end{equation}
so $d_A(H_{\mathrm{cpl}}) = 1$. The strict relation
\begin{equation}
d_A(H_{\mathrm{cpl}}) = 1 < 2 = d_A(H_{C_1}) + d_A(H_{C_2}) \tag{S28}
\end{equation}
is therefore caused by the emitted pair interaction together with the declared viable, resolution-limited quotient: it is not a consequence of changing from a two-coordinate to a one-coordinate readout by fiat. If the readout is coarsened below a resolvable phase step, the scale ladder cannot certify a unique covering exponent and the same construction reports $d_A = \bot$.

The stable invariant $q = \mathrm{sgn}\cos(\delta)$ in the phase-locked regime factors through the operational quotient. The finite presentation consisting of two clocks, the pair interaction node, the observer readout node, and their typed incidences has a faithful carrier in $\mathbb{R}^3$. Any phase-coordinate change preserving the coupled flow induces a dynamic-equivariant bridge and transports $q$ by Proposition 30. The witness therefore follows (S1) as one causal process: a global material action yields a composite interaction; the interaction changes the viable registration quotient; the quotient yields an invariant and a computed effective dimension; and the invariant can be represented and transported.

\noindent\plainterms{} The whole theory, run on the smallest interesting example. Two clocks on a wall: the extraction formula finds exactly one interaction (the wall's pull); the interaction locks the clocks together; locking makes the difference between them invisible to any patient observer at the stated precision, so the pair's tell-apart map shrinks from a two-dimensional grid to a one-dimensional loop --- the measured dimension drops from 2 to 1 because of the coupling, with nothing changed about the instruments; the locked-or-not fact is the invariant everyone agrees on; and the whole setup draws faithfully on paper. Coarsen the instruments too far and the construction says so honestly: unresolved.

\section{Boundary conditions and empirical realization}

Each condition here belongs to a particular production edge. The table records what a physical application must establish to instantiate that edge.

\begin{center}
\small
\begin{longtable}{@{}p{0.26\textwidth}p{0.34\textwidth}p{0.34\textwidth}@{}}
\toprule
Causal edge & Required condition & Application-level evidence\\
\midrule
\endhead
Local observations $\to$ global law & Declared diagram (L1)--(L5), including one grounding condition & Flow-compatible declared arrows, calibrated overlaps, a named grounding witness, material composition coherence\\
Global law $\to$ interactions & Generator regularity and localization & Partial-sum reconstruction, participant support, stable operational signatures\\
Flow $\to$ functional states & Monotone inclusive dependency and horizon labels & Closure test, observation-preserving transition matching\\
Functional states $\to$ HCFM & Material observer, sensitivity, viability, correction, canonical core construction & Protocol separation, readout quotient, concept-lattice and syntactic-category construction, continuation and calibration tests\\
Registration $\to d_A$ & Declared label cost, viable accessed region, scale ladder, tolerance, and covering-profile stability & A resolved value, an unresolved result, or a certified refinement-flow profile\\
Three-dimensional resolution & Saturation, exhaustion, and rank--metric regularity at rank 3 & Instantiation or blocking of Theorem 11.15; no inference if a premise is unverified\\
HCFM $\to$ disciplinary realization & Declared perspective, evaluator, equivalence, admissibility, and certificate interface & A type-audited realization under Definition 12.1\\
System hypothesis $\to$ coherence assessment & Declared HCFM system model, hypothesis target, perspective, coherence budget, noise floor, and observability ceiling & C, I, or UB under Definition 12.10\\
Recursive assessment $\to$ binary limit & Fairness, soundness, finite separation, completion, monotone convergence, and decisive resolution & Theorem 12.17, or a localized undetermined condition\\
Discipline $\to$ transport & A certified bridge satisfying (B1)--(B6) & Valid source-to-target scope and certificate record\\
Semantic library $\to$ certified solving-rate acceleration & C1--C6, nondegenerate modules, coherent task realization, activated descriptors, and update-work bound $q < 2$ & Theorem 12.27; measured throughput or a localized failed realization condition\\
Windows $\to$ transition or decoherence & Declared floor/ceiling, driven boundary encounter, independence audit of coupled structure & A nucleation-threshold estimate, a decoherence record, or refutation of Conjecture 13.10\\
Diffusion $\to$ net sign & Both rates of Definition 13.23 in bounded certified collectives & O1, O2, or O3, reported with equal prominence\\
HCFM $\to$ human carrier & Finite typed presentation & Decoding of the labelled $\mathbb{R}^3$ carrier; comparison by $J_H$ where desired\\
Invariant $\to$ transport & Typed dynamic-equivariant bridge & Commutation and invariant-preservation residuals\\
\bottomrule
\end{longtable}
\end{center}

\section{Final statement}

The formal construction establishes one navigable causal shape. Stable, compatible local material behavior assembles into one continuous global law. That law acts on each observable material composition and emits the interaction functions generated by composition. Closure and observational quotienting turn the lawful process into a functional state object. Physical observation turns that object into a composition-closed HCFM schema with canonical conceptual and mathematical cores, plus typed realized enrichment. Stable refinement carries the schema to a canonical global completion. A flat carrier exposes finite slices to human inspection, and structure-preserving bridges carry observable invariants to other systems.

The dimensionality module makes ``how many dimensions?'' a measurement inside that shape. Effective dimension is evaluated on a declared costed registration datum; the estimator-calibration theorem reports three when a rank-three certified separating family is rank--metric regular, while saturation and exhaustion fix the certified rank; the coupled-clock witness exhibits both the uncoupled product regime and a coupling-induced strict reduction of the measured value. On a declared realization class whose typical infrared trajectories satisfy those certificates, the prediction of a generic three-dimensional report is globally coherent; causal-order and ultraviolet refinement paths identify empirical realizations of that prediction. What was a metaphysical dispute is now a measurement protocol with named certificates and a positive conditional consequence.

The disciplinary module makes the shape cumulative, assessment-explicit, and blind-spot-aware. Disciplines are perspective-explicit functional realizations; cross-disciplinary reuse travels only over certified six-clause bridges that preserve evaluator, admissibility, semantic-equivalence, and certificate conditions. The observer-inclusion blind spot is the systemic assessment failure mode in which those conditions remain outside an HCFM system model: a one-shot local-prior judgment can then require an unbounded sequence of certified corrections before the model and assessment perspective are jointly explicit. The coherence-budgeted assessment protocol returns C, I, or UB for a hypothesis about an internally executed or externally declared HCFM system model; external certification serializes the same constraint algebra for independent audit. Its blind-spot proposition gives sufficient conditions for unbounded discrepancy; its recursive global-coherence convergence theorem identifies the fairness, soundness, finite-separation, completion, convergence, and resolution conditions under which a binary verdict is reached. The semantic-library mechanism is a bounded, deduplicated, certified path-composition system whose best retained answers provably never degrade, whose certified candidate capacity is exponential under six named conditions, and whose certified real-task throughput accelerates under the coherent task-realization and update-work conditions of Theorem 12.27. The dependency records and realization ladder let narrower papers be projected from this reference with their inferential ancestry intact while preserving the positive rate prediction and its empirical measurement terms.

The alignment module turns the same machinery on the growth of capable systems. It defines operational windows on realized horizons, types the boundary alternatives as coherent meta-lift or loss of the operational quotient, and proves that cognitive mass --- signed reach change per unit of registered coupling --- is an order-relative observable invariant with a two-channel population decomposition whose realization term is the discontinuous one. On that proved base it registers, with falsifiers and dependency chains, the nucleation requirement, the two-lever control claim, the protocol-internal realization gate, the anti-foom corollary, the survivorship reading of the Great Filter, and the preregistered sign question on which the priority of the two levers turns.

The result is a conditional Theory-of-Everything construction in which one law and one schema generate the physical, observational, representational, transport, dimensional, disciplinary, and alignment-dynamical claims --- with the type of each claim locally explicit, the ledgers holding the bookkeeping, and every arrow of the completed architecture pointing outward at a named test.

\begin{thebibliography}{99}
\bibitem{1} A. Tarski. A lattice-theoretical fixpoint theorem and its applications. \emph{Pacific Journal of Mathematics}, 5:285--309, 1955.
\bibitem{2} P. Ehrenfest. In what way does it become manifest in the fundamental laws of physics that space has three dimensions? \emph{Proceedings of the Amsterdam Academy}, 20:200--209, 1917.
\bibitem{3} G. J. Whitrow. Why physical space has three dimensions. \emph{British Journal for the Philosophy of Science}, 6(21):13--31, 1955.
\bibitem{4} M. Tegmark. On the dimensionality of spacetime. \emph{Classical and Quantum Gravity}, 14(4):L69--L75, 1997.
\bibitem{5} N. Arkani-Hamed, S. Dimopoulos, and G. Dvali. The hierarchy problem and new dimensions at a millimeter. \emph{Physics Letters B}, 429(3--4):263--272, 1998.
\bibitem{6} L. Randall and R. Sundrum. Large mass hierarchy from a small extra dimension. \emph{Physical Review Letters}, 83(17):3370--3373, 1999.
\bibitem{7} J. Ambj\o rn, J. Jurkiewicz, and R. Loll. Spectral dimension of the universe. \emph{Physical Review Letters}, 95(17):171301, 2005.
\bibitem{8} O. Lauscher and M. Reuter. Fractal spacetime structure in asymptotically safe gravity. \emph{Journal of High Energy Physics}, 2005(10):050, 2005.
\bibitem{9} S. Carlip. Dimension and dimensional reduction in quantum gravity. \emph{Classical and Quantum Gravity}, 34(19):193001, 2017.
\bibitem{10} J. Myrheim. Statistical geometry. CERN preprint TH-2538, 1978.
\bibitem{11} D. A. Meyer. \emph{The Dimension of Causal Sets}. PhD thesis, Massachusetts Institute of Technology, 1988.
\bibitem{12} L. Bombelli, J. Lee, D. Meyer, and R. D. Sorkin. Space-time as a causal set. \emph{Physical Review Letters}, 59(5):521--524, 1987.
\bibitem{13} A. E. Williams. Cognitive Near-Singularity: A possibility theorem and a witness-based protocol with a corpus case study. \emph{Frontiers in Artificial Intelligence}, 2026.
\bibitem{14} H. von Foerster. \emph{Cybernetics of Cybernetics}. Biological Computer Laboratory Report 73.38, University of Illinois, Urbana, IL, 1974.
\bibitem{15} J. Pearl. Causal inference in statistics: An overview. \emph{Statistics Surveys}, 3:96--146, 2009. doi:10.1214/09-SS057.
\bibitem{16} C. F. Manski. \emph{Partial Identification of Probability Distributions}. Springer-Verlag, New York, 2003.
\bibitem{17} A. Brandenburger and E. Dekel. Hierarchies of beliefs and common knowledge. \emph{Journal of Economic Theory}, 59(1):189--198, 1993. doi:10.1006/jeth.1993.1012.
\bibitem{18} E. M. Gold. Language identification in the limit. \emph{Information and Control}, 10(5):447--474, 1967. doi:10.1016/S0019-9958(67)91165-5.
\bibitem{19} H. B. Enderton. \emph{A Mathematical Introduction to Logic}. 2nd edition, Academic Press, San Diego, CA, 2001.
\bibitem{20} P. Cousot and R. Cousot. Abstract interpretation: A unified lattice model for static analysis of programs by construction or approximation of fixpoints. In \emph{Proceedings of the 4th ACM SIGACT--SIGPLAN Symposium on Principles of Programming Languages}, pages 238--252, 1977.
\bibitem{21} A. Finkel and P. Schnoebelen. Well-structured transition systems everywhere! \emph{Theoretical Computer Science}, 256(1--2):63--92, 2001.
\bibitem{22} L. von Bertalanffy. \emph{General System Theory: Foundations, Development, Applications}. George Braziller, New York, 1968.
\bibitem{23} B. Fong and D. I. Spivak. \emph{An Invitation to Applied Category Theory: Seven Sketches in Compositionality}. Cambridge University Press, 2019.
\bibitem{24} J. C. Willems. The behavioral approach to open and interconnected systems. \emph{IEEE Control Systems Magazine}, 27(6):46--99, 2007.
\bibitem{25} N. Bourbaki. \emph{General Topology, Chapters 1--4}. Springer, Berlin, 1995.
\bibitem{26} S. Eilenberg and N. Steenrod. \emph{Foundations of Algebraic Topology}. Princeton University Press, 1952.
\bibitem{27} L. Ribes and P. Zalesskii. \emph{Profinite Groups}. 2nd edition, Springer, Berlin, 2010.
\bibitem{28} L. Henkin. A problem on inverse mapping systems. \emph{Proceedings of the American Mathematical Society}, 1:224--225, 1950.
\bibitem{29} W. C. Waterhouse. An empty inverse limit. \emph{Proceedings of the American Mathematical Society}, 36(2):618, 1972.
\bibitem{30} B. Ganter and R. Wille. \emph{Formal Concept Analysis: Mathematical Foundations}. Springer, Berlin, 1999.
\bibitem{31} J. Lambek and P. J. Scott. \emph{Introduction to Higher Order Categorical Logic}. Cambridge University Press, 1986.
\bibitem{32} P. T. Johnstone. \emph{Sketches of an Elephant: A Topos Theory Compendium, Vol. 2}. Oxford University Press, 2002.
\bibitem{33} Y. Kuramoto. \emph{Chemical Oscillations, Waves, and Turbulence}. Springer, Berlin, 1984.
\bibitem{34} A. Pikovsky, M. Rosenblum, and J. Kurths. \emph{Synchronization: A Universal Concept in Nonlinear Sciences}. Cambridge University Press, 2001.
\bibitem{35} S. H. Strogatz. From Kuramoto to Crawford: exploring the onset of synchronization in populations of coupled oscillators. \emph{Physica D}, 143(1--4):1--20, 2000.
\bibitem{36} J. A. Acebr\'on, L. L. Bonilla, C. J. P\'erez Vicente, F. Ritort, and R. Spigler. The Kuramoto model: a simple paradigm for synchronization phenomena. \emph{Reviews of Modern Physics}, 77(1):137--185, 2005.
\end{thebibliography}

\end{document}
